% Lean desde cero: de #check a los teoremas de la tesis. Formato visual,
% amigable para TDAH e interactivo (PDF con enlaces). 19 niveles en 4 mundos.
% El código es el de Thesis/Prop y Thesis/Sort; los estados de prueba son los
% que muestra Lean (exportados en artifacts/explorer o comprobados aparte).
% Los bloques lean/estado se colorean solos con un filtro Lua.
% Compilar dos veces:  lualatex LeanDesdeCero.tex
\documentclass[12pt]{article}
\usepackage[a4paper,landscape,margin=11mm,bottom=17mm]{geometry}
\usepackage{amsmath}
\usepackage{fontspec}
\usepackage{unicode-math}
\usepackage[spanish,es-noshorthands]{babel}
\directlua{
  luaotfload.add_fallback("leanfallback", {
    "STIX Two Math:mode=harf;",
    "DejaVu Sans:mode=harf;",
    "Apple Symbols:mode=harf;",
  })
}
\setmainfont{TeX Gyre Heros}[RawFeature={fallback=leanfallback}]
\setsansfont{TeX Gyre Heros}[RawFeature={fallback=leanfallback}]
\setmathfont{TeX Gyre DejaVu Math}
\setmonofont{DejaVu Sans Mono}[Scale=MatchLowercase, RawFeature={fallback=leanfallback}]
\usepackage[table]{xcolor}
\usepackage{fvextra}
\usepackage{luacode}
\usepackage[most]{tcolorbox}
\usepackage{tikz}
\usetikzlibrary{arrows.meta,positioning,calc,shapes.geometric,trees}
\usepackage{fontawesome5}
\usepackage{enumitem}
\usepackage{array}
\usepackage{graphicx}
\usepackage{bussproofs}
\usepackage{lastpage}
\usepackage[hidelinks]{hyperref}
\pagestyle{empty}
\setlength{\parindent}{0pt}
\linespread{1.06}

% ---------- Colores: uno por mundo ----------
\definecolor{m1}{HTML}{3F6FA8}   % Mundo 1: leer Lean
\definecolor{m2}{HTML}{8C5FB8}   % Mundo 2: probar
\definecolor{m3}{HTML}{2A8C8C}   % Mundo 3: inducción y recursión
\definecolor{m4a}{HTML}{3E9A5A}  % Mundo 4: completitud
\definecolor{m4b}{HTML}{D9822B}  % Mundo 4: quicksort
\definecolor{m4c}{HTML}{C8463D}  % Mundo 4: auditoría
\definecolor{gris}{HTML}{5B6B82}
\definecolor{tldr}{HTML}{FFF4C2}
\definecolor{verde}{HTML}{3E9A5A}
\definecolor{ambar}{HTML}{E0A320}
\definecolor{rojo}{HTML}{C8463D}
\definecolor{kw}{HTML}{1F4E9C}
\definecolor{tac}{HTML}{8A2B8F}
\definecolor{com}{HTML}{5F6F5F}
\definecolor{fondocodigo}{HTML}{F6F7FA}
\definecolor{bordecodigo}{HTML}{9AA7B8}
\definecolor{fondoestado}{HTML}{FFF8E6}
\definecolor{bordeestado}{HTML}{D9B45A}
\foreach \i in {1,...,5}{\xglobal\colorlet{n\i}{m1}}
\foreach \i in {6,...,10}{\xglobal\colorlet{n\i}{m2}}
\foreach \i in {11,...,14}{\xglobal\colorlet{n\i}{m3}}
\foreach \i in {15,16}{\xglobal\colorlet{n\i}{m4a}}
\foreach \i in {17,18}{\xglobal\colorlet{n\i}{m4b}}
\colorlet{n19}{m4c}

% ---------- Coloreado de código Lean (filtro Lua sobre las líneas) ----------
\newcommand{\LK}[1]{\textcolor{kw}{\bfseries #1}}
\newcommand{\LT}[1]{\textcolor{tac}{\bfseries #1}}
\newcommand{\LC}[1]{\textcolor{com}{#1}}
\newcommand{\LG}[1]{\textcolor{rojo!85!black}{\bfseries #1}}
\newcommand{\LH}[1]{\colorbox{tldr}{#1}}
\newcommand{\LB}[1]{\tikz[baseline=(n.base)]\node[circle, fill=black!78, text=white,
  inner sep=0.8pt, minimum size=3.6mm, font=\sffamily\bfseries\scriptsize](n){#1};}
\newcommand{\LN}[1]{{\normalfont\sffamily\footnotesize\color{gris}\ \faArrowLeft\ #1}}
\newcommand{\bola}[1]{\LB{#1}}
\begin{luacode*}
local envs = { leanL = "code", lean = "code", leanS = "code", leanXS = "code",
               estadoL = "state", estado = "state", estadoS = "state" }
local mode, cur = nil, nil
local KW = {}
for w in ([[theorem lemma def abbrev inductive structure where by with fun match
  from at in if then else termination_by decreasing_by namespace open variable
  end deriving notation infixr infixl infix prefix noncomputable import section
  example omit attribute instance]]):gmatch("%S+") do KW[w] = "LK" end
for w in ([[intro intros exact apply refine constructor cases rcases obtain rintro
  induction rw simp simpa change subst by_cases by_contra rfl have let show calc
  sorry]]):gmatch("%S+") do KW[w] = "LT" end
local function esc(s)
  return (s:gsub("[\\{}]", { ["\\"] = "\\textbackslash{}", ["{"] = "\\{", ["}"] = "\\}" }))
end
local function hlcode(s)
  local out = s:gsub("()([%a_][%w_']*)", function(pos, w)
    if pos > 1 and s:sub(pos - 1, pos - 1) == "." then return nil end
    local k = KW[w]
    if k then return "\1" .. k .. "\2" .. w .. "\3" end
  end)
  out = out:gsub("#(%a+)", "\1LK\2#%1\3")
  out = esc(out)
  out = out:gsub("\1", "\\"):gsub("\2", "{"):gsub("\3", "}")
  return out
end
local function codeline(body)
  if body:match("^%s*/%-") then return "\\LC{" .. esc(body) .. "}" end
  local i = body:find("--", 1, true)
  if i then
    return hlcode(body:sub(1, i - 1)) .. "\\LC{" .. esc(body:sub(i)) .. "}"
  end
  return hlcode(body)
end
local function stateline(body)
  local e = esc(body)
  if body:match("^%s*⊢") then return "\\LG{" .. e .. "}" end
  if body:match("^case ") then return "\\LC{" .. e .. "}" end
  return e
end
local function filter(line)
  if mode then
    if line:match("^%s*\\end{" .. cur .. "}") then mode = nil; return line end
    local note, ball, hi = "", "", false
    local b, n = line:match("^(.-)%s*@@%s*(.*)$")
    if b then line = b; note = "\\LN{" .. esc(n) .. "}" end
    b, n = line:match("^(.-)@(%d+)%s*$")
    if b then line = b; ball = "\\LB{" .. n .. "}" end
    b = line:match("^(.-)%s*@!%s*$")
    if b then line = b; hi = true end
    local out = (mode == "code") and codeline(line) or stateline(line)
    if hi then out = "\\LH{" .. out .. "}" end
    return out .. ball .. note
  else
    local name = line:match("^%s*\\begin{(%a+)}")
    if name and envs[name] then mode = envs[name]; cur = name end
    return line
  end
end
luatexbase.add_to_callback("process_input_buffer", filter, "leanhl")
\end{luacode*}

\tcbset{
  leanbox/.style={enhanced, colback=fondocodigo, colframe=bordecodigo, boxrule=0.9pt,
    arc=4pt, left=5pt, right=3pt, top=-2pt, bottom=-2pt, before skip=3pt, after skip=4pt,
    fonttitle=\ttfamily\footnotesize, coltitle=black!75, colbacktitle=bordecodigo!30,
    toptitle=1pt, bottomtitle=1pt},
  estadobox/.style={enhanced, colback=fondoestado, colframe=bordeestado, boxrule=0.9pt,
    arc=4pt, left=5pt, right=3pt, top=-2pt, bottom=-2pt, before skip=3pt, after skip=4pt,
    fonttitle=\sffamily\bfseries\footnotesize, coltitle=black!80,
    colbacktitle=bordeestado!45, toptitle=1pt, bottomtitle=1pt},
}
\newcommand{\defleanenv}[3]{%
  \newenvironment{#1}[1][]{\VerbatimEnvironment
    \begin{tcolorbox}[#3,##1]\begin{Verbatim}[commandchars=\\\{\},fontsize=#2]}
   {\end{Verbatim}\end{tcolorbox}}}
\defleanenv{leanL}{\fontsize{11.5}{14.2}\selectfont}{leanbox}
\defleanenv{lean}{\fontsize{10.5}{13}\selectfont}{leanbox}
\defleanenv{leanS}{\fontsize{9.6}{11.9}\selectfont}{leanbox}
\defleanenv{leanXS}{\fontsize{8.8}{10.9}\selectfont}{leanbox}
\defleanenv{estadoL}{\fontsize{11}{13.6}\selectfont}{estadobox}
\defleanenv{estado}{\fontsize{10}{12.4}\selectfont}{estadobox}
\defleanenv{estadoS}{\fontsize{9.4}{11.6}\selectfont}{estadobox}

\newcommand{\cod}[1]{\texttt{\detokenize{#1}}}
\AtBeginDocument{\providecommand\square{\mdlgwhtsquare}}

% ---------- Botones y navegación ----------
\newcommand{\boton}[2][gris]{\tikz[baseline=(b.base)]\node[draw=#1, fill=#1!10,
  line width=1pt, rounded corners=5pt, inner xsep=8pt, inner ysep=4pt,
  font=\bfseries\normalsize, text=#1!80!black](b){#2};}
\newcommand{\ir}[3][gris]{\hyperlink{#2}{\boton[#1]{#3}}}
\newcommand{\chip}[1]{\hyperlink{nivel#1}{\tikz[baseline=(c.base)]\node[fill=n#1,
  text=white, rounded corners=3pt, inner sep=2.5pt, font=\bfseries\footnotesize](c){#1};}}
\makeatletter\newcommand\pdim[1]{\strip@pt\dimexpr#1\relax}\makeatother
\newcommand{\navegacion}[4]{\AddToHookNext{shipout/foreground}{%
  \put(\pdim{0.5\paperwidth},\pdim{-\paperheight+4mm}){\makebox(0,0)[b]{%
  \ir{#1}{\faArrowLeft\ #2}\hspace{5mm}\ir{mapa}{\faMapSigns\ Mapa}%
  \hspace{5mm}\ir{simbolos}{\faKeyboard\ Símbolos}\hspace{5mm}\ir{diccionario}{\faCommentDots\ Cómo decirlo}%
  \hspace{5mm}\ir{indice}{\faFont\ Índice}\hspace{5mm}\ir{#3}{#4\ \faArrowRight}}}}}

% ---------- Encabezados ----------
% #1 número, #2 mundo, #3 icono, #4 título, #5 para qué sirve
\newcommand{\nivel}[5]{%
  \begin{tcolorbox}[colback=n#1, colframe=n#1, arc=5pt, boxrule=0pt,
      left=10pt, right=10pt, top=4pt, bottom=4pt, after skip=3pt]
    \hypertarget{nivel#1}{}\color{white}\LARGE #3\quad\bfseries #1.\ #4\hfill
    \normalfont\large\faGamepad\ nivel #1 de 19 · mundo #2
  \end{tcolorbox}
  \begin{tcolorbox}[colback=tldr, colframe=n#1, boxrule=0.8pt, arc=4pt,
      left=10pt, right=10pt, top=3pt, bottom=3pt, after skip=5pt]
    \large\faBullseye\ \textbf{Para qué sirve:}\ #5
  \end{tcolorbox}}
\newcommand{\pagina}[5]{%
  \begin{tcolorbox}[colback=#2, colframe=#2, arc=5pt, boxrule=0pt,
      left=10pt, right=10pt, top=4pt, bottom=4pt, after skip=3pt]
    \hypertarget{#1}{}\color{white}\LARGE #3\quad\bfseries #4\hfill
    \normalfont\large página \thepage\ de \pageref{LastPage}
  \end{tcolorbox}
  \begin{tcolorbox}[colback=tldr, colframe=#2, boxrule=0.8pt, arc=4pt,
      left=10pt, right=10pt, top=3pt, bottom=3pt, after skip=5pt]
    \large\faBolt\ \textbf{En 10 segundos:}\ #5
  \end{tcolorbox}}

% ---------- Piezas de contenido ----------
% Explicación numerada: \ex{n}{texto}
\newcommand{\ex}[2]{\par\noindent\hangindent=7mm\hangafter=1\makebox[7mm][l]{\bola{#1}}#2\par\vspace{2.5pt}}
% Tarjeta: \tarj{color}{título}{texto}
\newcommand{\tarj}[3]{\begin{tcolorbox}[colback=#1!6, colframe=#1, boxrule=1pt,
  arc=4pt, left=5pt, right=5pt, top=3pt, bottom=3pt, leftrule=5pt,
  before upper=\raggedright, before skip=3pt, after skip=4pt]
  {\bfseries\color{#1!80!black}#2}\par\vspace{1pt}#3\end{tcolorbox}}
% Idea clave
\newcommand{\clave}[2]{\begin{tcolorbox}[colback=tldr, colframe=#1, boxrule=1pt, arc=4pt,
  left=6pt, right=6pt, top=3pt, bottom=3pt, before skip=3pt, after skip=4pt,
  before upper=\raggedright]\faLightbulb\ #2\end{tcolorbox}}
% Cómo se dice
\newcommand{\dice}[1]{\begin{tcolorbox}[colback=white, colframe=gris, boxrule=1pt, arc=4pt,
  left=6pt, right=6pt, top=3pt, bottom=3pt, before skip=3pt, after skip=4pt,
  before upper=\raggedright]{\color{gris}\faCommentDots}\ \textbf{Se dice:} #1\end{tcolorbox}}
% Pruébate
\newcommand{\pruebate}[3]{%
  \begin{tcolorbox}[colback=white, colframe=#1, boxrule=1.1pt, arc=5pt,
    left=8pt, right=8pt, top=3pt, bottom=3pt, sidebyside, lefthand width=0.56\linewidth,
    sidebyside gap=10pt, segmentation style={solid, #1!40}, before skip=4pt]
    {\large\bfseries\color{#1!80!black}\faBrain\ Pruébate} \ {\small(dilo en voz alta)}
    \begin{enumerate}[leftmargin=6mm, itemsep=1pt, topsep=2pt]#2\end{enumerate}
    \tcblower
    \centering{\footnotesize\color{black!50}\faSync*\ respuestas: gira la página}\par\vspace{2pt}
    \rotatebox{180}{\parbox{0.92\linewidth}{\small\raggedright #3}}
  \end{tcolorbox}}
\tikzset{num/.style={circle, fill=#1, text=white, font=\bfseries\Large,
    minimum size=9mm, inner sep=0pt}}

\begin{document}

%=====================================================================
% PORTADA
%=====================================================================
\hypertarget{portada}{}
\begin{tcolorbox}[colback=m1, colframe=m1, arc=8pt, boxrule=0pt, halign=center,
  top=12pt, bottom=12pt]
\color{white}{\fontsize{34}{40}\selectfont\bfseries\faCode\ Lean desde cero}\\[8pt]
{\LARGE De \texttt{\#check} a los teoremas de la tesis}\\[4pt]
{\large 19 niveles · 4 mundos · código real de la tesis · todo es clicable}
\end{tcolorbox}

\vspace{4pt}
\begin{minipage}[t]{0.50\textwidth}
{\Large\bfseries Cómo se usa}\\[6pt]
\begin{tikzpicture}[baseline=(current bounding box.north),
    txt/.style={align=left, text width=11.4cm, anchor=west, font=\large}]
  \node[num=gris] (a) at (0,0) {1};
  \node[txt, right=4mm of a] {Abre \textbf{el mapa} y juega los niveles en orden.\\ Cada nivel usa solo lo de los anteriores.};
  \node[num=gris] (b) at (0,-1.5) {2};
  \node[txt, right=4mm of b] {En cada nivel: \textbf{para qué sirve} $\to$ \textbf{código real}\\ $\to$ \textbf{estado de Lean} $\to$ \textbf{pruébate}.};
  \node[num=gris] (c) at (0,-3.0) {3};
  \node[txt, right=4mm of c] {Los números \bola{1} \bola{2} \bola{3} del código tienen su\\ explicación al lado, con el mismo número.};
  \node[num=gris] (d) at (0,-4.5) {4};
  \node[txt, right=4mm of d] {¿Poco tiempo? Sigue la \textbf{ruta rápida} del mapa.};
\end{tikzpicture}
\end{minipage}\hfill
\begin{minipage}[t]{0.46\textwidth}
{\Large\bfseries Empieza por aquí (haz clic)}\\[8pt]
\begin{tcolorbox}[colback=m1, colframe=m1, arc=6pt, halign=center, hyperlink=mapa]
  \color{white}\Large\faMapSigns\ \textbf{El mapa de niveles}\\ \normalsize 19 niveles, de leer a demostrar\end{tcolorbox}
\begin{tcbraster}[raster columns=3, raster column skip=6pt, raster equal height]
\begin{tcolorbox}[colback=m2, colframe=m2, arc=6pt, halign=center, hyperlink=simbolos]
  \color{white}\large\faKeyboard\ \textbf{Símbolos}\\ \small cómo se leen\end{tcolorbox}
\begin{tcolorbox}[colback=m3, colframe=m3, arc=6pt, halign=center, hyperlink=diccionario]
  \color{white}\large\faCommentDots\ \textbf{Cómo decirlo}\\ \small y 10 trampas\end{tcolorbox}
\begin{tcolorbox}[colback=m4c, colframe=m4c, arc=6pt, halign=center, hyperlink=indice]
  \color{white}\large\faFont\ \textbf{Índice}\\ \small táctica $\to$ nivel\end{tcolorbox}
\end{tcbraster}

\vspace{2pt}
{\large\bfseries Las tres cajas que verás}\\[2pt]
\begin{tcbraster}[raster columns=3, raster column skip=6pt, raster equal height]
\begin{tcolorbox}[leanbox, top=3pt, bottom=3pt, halign=center]\small\textbf{gris}: código Lean\end{tcolorbox}
\begin{tcolorbox}[estadobox, top=3pt, bottom=3pt, halign=center]\small\textbf{amarilla}: lo que responde Lean\end{tcolorbox}
\begin{tcolorbox}[colback=white, colframe=gris, boxrule=1pt, arc=4pt, top=3pt, bottom=3pt, halign=center]\small{\color{gris}\faCommentDots}\ cómo se dice\end{tcolorbox}
\end{tcbraster}
\end{minipage}

\vfill
\begin{tcolorbox}[colback=tldr, colframe=gris, arc=4pt, boxrule=0.8pt, halign=center]
\large\faStar\ \textbf{Regla de oro:} para explicar una prueba de Lean, describe el
\textbf{estado} (qué tienes y qué falta), no la táctica.
\end{tcolorbox}

\newpage
%=====================================================================
% MAPA
%=====================================================================
\pagina{mapa}{gris}{\faMapSigns}{El mapa: 19 niveles en 4 mundos}{haz clic en un
nivel. Primero aprendes a leer, luego a probar, luego la inducción, y al final lees los teoremas.}

% #1 x, #2 y, #3 número, #4 icono, #5 título
\newcommand{\nd}[5]{%
  \node[fill=n#3, rounded corners=6pt, minimum width=4.85cm, minimum height=1.95cm,
        inner sep=2pt] (L#3) at (#1,#2)
   {\hyperlink{nivel#3}{\parbox[c][1.75cm][c]{4.65cm}{\centering\color{white}%
      {\Large #4\ \ \bfseries #3}\\[2pt]{\large\bfseries #5}}}};}
\begin{center}
\begin{tikzpicture}
  \node[anchor=west, font=\large\bfseries, text=m1] at (-2.4,1.35) {\faBookOpen\ MUNDO 1 · Leer Lean};
  \nd{0}{0}{1}{\faLightbulb}{La idea: término : tipo}
  \nd{5.45}{0}{2}{\faFile*}{Anatomía de un archivo}
  \nd{10.9}{0}{3}{\faCubes}{Datos: inductive}
  \nd{16.35}{0}{4}{\faCogs}{Funciones: def}
  \nd{21.8}{0}{5}{\faQuoteRight}{Enunciar: theorem}
  \node[anchor=west, font=\large\bfseries, text=m2] at (-2.4,-1.95) {\faHammer\ MUNDO 2 · Probar};
  \nd{0}{-3.3}{6}{\faEye}{El estado de prueba}
  \nd{5.45}{-3.3}{7}{\faHandPointRight}{Cinco tácticas}
  \nd{10.9}{-3.3}{8}{\faTable}{La tabla maestra}
  \nd{16.35}{-3.3}{9}{\faEquals}{Reescribir y calcular}
  \nd{21.8}{-3.3}{10}{\faCodeBranch}{Casos y absurdo}
  \node[anchor=west, font=\large\bfseries, text=m3] at (-2.4,-5.25) {\faRedo\ MUNDO 3 · Inducción y recursión};
  \nd{0}{-6.6}{11}{\faListOl}{Inducción sobre datos}
  \nd{5.45}{-6.6}{12}{\faProjectDiagram}{Inducción sobre reglas}
  \nd{10.9}{-6.6}{13}{\faSync*}{Recursión bien fundada}
  \nd{16.35}{-6.6}{14}{\faTree}{Construir derivaciones}
  \node[anchor=west, font=\large\bfseries, text=m4a] at (-2.4,-8.55) {\faTrophy\ MUNDO 4 · Los teoremas de la tesis};
  \nd{0}{-9.9}{15}{\faPuzzlePiece}{Completitud: las piezas}
  \nd{5.45}{-9.9}{16}{\faCheckDouble}{Completitud: el teorema}
  \nd{10.9}{-9.9}{17}{\faExchange*}{Quicksort: permutación}
  \nd{16.35}{-9.9}{18}{\faSortAmountUp}{Quicksort: orden y final}
  \nd{21.8}{-9.9}{19}{\faShield*}{Auditoría}
  \foreach \a/\b in {1/2,2/3,3/4,4/5,6/7,7/8,8/9,9/10,11/12,12/13,13/14,15/16,16/17,17/18,18/19}
    \draw[-{Stealth[length=3mm]}, line width=1.6pt, black!40] (L\a) -- (L\b);
\end{tikzpicture}
\end{center}

\vspace{-2pt}
\begin{tcolorbox}[colback=tldr, colframe=gris, arc=4pt, boxrule=0.8pt, halign=center]
\large\faBolt\ \textbf{Ruta rápida (20 minutos):}\quad
\chip{1} $\to$ \chip{5} $\to$ \chip{6} $\to$ \chip{8} $\to$ \chip{12} $\to$ \chip{16} $\to$ \chip{17} $\to$ \chip{19}
\end{tcolorbox}

\navegacion{portada}{Portada}{nivel1}{Nivel 1}

\newpage
%=====================================================================
% NIVEL 1
%=====================================================================
\nivel{1}{1}{\faLightbulb}{La idea: todo es «término : tipo»}{es la única idea de Lean.
Si entiendes \cod{a : A}, todo lo demás son variaciones.}

\begin{minipage}[t]{0.47\textwidth}
\vspace{0pt}
\begin{tcolorbox}[colback=white, colframe=black!75, boxrule=1.4pt, arc=6pt, halign=center,
  top=4pt, bottom=4pt]
{\fontsize{30}{34}\selectfont\ttfamily a : A}\\[2pt]
\large «\textbf{a} tiene tipo \textbf{A}»
\end{tcolorbox}
\begin{leanL}[title={Pregúntale a Lean con \#check y \#eval}]
#check 2 + 2                  -- 2 + 2 : ℕ
#check Nat                    -- Nat : Type
#check 2 + 2 = 4              -- 2 + 2 = 4 : Prop
#eval  2 + 2                  -- 4
#check Formula.atom "P"       -- … : Formula
#check P ⟶ P                  -- P ⟶ P : Formula
#check EsTautologia (P ⟶ P)   -- … : Prop
\end{leanL}
\tarj{m1}{\texttt{\#check e}}{Lean responde el \textbf{tipo} de \cod{e}. No calcula nada.}
\tarj{m1}{\texttt{\#eval e}}{Lean \textbf{calcula} el valor de \cod{e}.}
\end{minipage}\hfill
\begin{minipage}[t]{0.50\textwidth}
\vspace{0pt}
\begin{center}
\begin{tikzpicture}[
  c/.style={draw=#1, fill=#1!8, line width=1.2pt, rounded corners=5pt, align=center,
     text width=5.6cm, minimum height=1.5cm, inner sep=4pt},
  up/.style={-{Stealth[length=3mm]}, line width=1.5pt, black!55}]
  \node[font=\Large\bfseries, text=m1] at (0,4.55) {DATOS};
  \node[font=\Large\bfseries, text=m2] at (6.6,4.55) {PROPOSICIONES};
  \node[c=m1] (T) at (0,3.3) {{\Large\ttfamily Type}\\ el tipo de los tipos de datos};
  \node[c=m2] (P) at (6.6,3.3) {{\Large\ttfamily Prop}\\ el tipo de las proposiciones};
  \node[c=m1] (t1) at (0,0.95) {\cod{Nat} \ \ \cod{Formula} \ \ \cod{List α}\\ tipos de datos};
  \node[c=m2] (p1) at (6.6,0.95) {\cod{2 + 2 = 4} \ \ \cod{ND Γ φ}\\ enunciados};
  \node[c=m1] (t0) at (0,-1.4) {\cod{4} \ \ \cod{P ⟶ P} \ \ \cod{[3, 1, 2]}\\ \textbf{datos}};
  \node[c=m2] (p0) at (6.6,-1.4) {\cod{rfl} \ \ \cod{completitud_ND φ h}\\ \textbf{pruebas}};
  \draw[up] (t1) -- node[right, font=\ttfamily\large]{:} (T);
  \draw[up] (p1) -- node[right, font=\ttfamily\large]{:} (P);
  \draw[up] (t0) -- node[right, font=\ttfamily\large]{:} (t1);
  \draw[up] (p0) -- node[right, font=\ttfamily\large]{:} (p1);
\end{tikzpicture}
\end{center}
\clave{m1}{\large\textbf{Una prueba es un término cuyo tipo es el enunciado.}
Demostrar $=$ construir ese término. Lean solo comprueba que el tipo encaja.}
\end{minipage}

\pruebate{m1}{\item ¿Qué tipo tiene \cod{P ⟶ P}? ¿Y \cod{EsTautologia (P ⟶ P)}?
\item ¿Qué diferencia hay entre \texttt{\#check} y \texttt{\#eval}?
\item ¿Qué es «demostrar» para Lean?}
{1.\ \cod{Formula} (es un dato); \cod{Prop} (es un enunciado).
2.\ \texttt{\#check} da el tipo; \texttt{\#eval} calcula el valor.
3.\ Dar un término cuyo tipo es el enunciado.}

\navegacion{mapa}{Mapa}{nivel2}{Nivel 2}

\newpage
%=====================================================================
% NIVEL 2
%=====================================================================
\nivel{2}{1}{\faFile*}{Anatomía de un archivo}{reconocer las partes de un archivo
\texttt{.lean} para no perderte al abrir el código de la tesis.}

\begin{minipage}[t]{0.47\textwidth}
\vspace{0pt}
\begin{leanL}[title={Thesis/Prop/Syntax.lean (el archivo completo)}]
namespace Thesis.Prop                           @1

inductive Formula : Type                        @2
| atom : String → Formula
| neg  : Formula → Formula
| impl : Formula → Formula → Formula
deriving DecidableEq, Repr                      @3

prefix:60 "~" => Formula.neg                    @4
infixr:55 " ⟶ " => Formula.impl

abbrev Valuacion : Type := String → Prop        @5

def evaluar (v : Valuacion) : Formula → Prop    @6
| .atom s   => v s
| .neg p    => ¬ evaluar v p
| .impl p q => evaluar v p → evaluar v q

def EsTautologia (φ : Formula) : Prop :=        @7
  ∀ v : Valuacion, evaluar v φ

end Thesis.Prop
\end{leanL}
\end{minipage}\hfill
\begin{minipage}[t]{0.50\textwidth}
\vspace{0pt}
\ex{1}{\cod{namespace … end}: todo lo de dentro se llama \cod{Thesis.Prop.…}. Con \cod{open} se omite el prefijo.}
\ex{2}{\cod{inductive}: un tipo de datos nuevo, con tres constructores \chip{3}.}
\ex{3}{\cod{deriving}: Lean genera solo la igualdad decidible y cómo imprimirlo.}
\ex{4}{\textbf{Notación}: \cod{~φ} y \cod{φ ⟶ ψ} abrevian \cod{Formula.neg} y \cod{Formula.impl}. El 60 y el 55 son la prioridad: \cod{~P ⟶ Q} es \cod{(~P) ⟶ Q}.}
\ex{5}{\cod{abbrev}: un alias. Una valuación es una función de nombres de átomo en \cod{Prop}.}
\ex{6}{\cod{def} con casos: una función recursiva \chip{4}.}
\ex{7}{Una definición cuyo valor es una \textbf{proposición}.}
\vspace{2pt}
\tarj{m1}{Lo que no sale aquí}{\cod{import}: trae otro archivo (siempre arriba). \cod{-- texto}: comentario.
\cod{/-- texto -/}: documentación. \cod{-- ANCHOR:}: marcas para citar el código en la tesis.}
\begin{tcolorbox}[colback=m1!6, colframe=m1, boxrule=1pt, arc=4pt, left=6pt, right=6pt,
  top=3pt, bottom=3pt, title={\faFolderOpen\ Los archivos de la tesis}, colbacktitle=m1,
  fonttitle=\bfseries]
\small
\cod{Thesis/Prop/Syntax.lean} \hfill fórmulas y semántica\\
\cod{Thesis/Prop/NaturalDeduction.lean} \hfill el cálculo \cod{ND} y la corrección\\
\cod{Thesis/Prop/Completeness.lean} \hfill Kalmár, descarga, completitud\\
\cod{Thesis/Prop/Main.lean} \hfill ejemplos \quad \cod{Thesis/Prop/Audit.lean} \hfill axiomas\\
\cod{Thesis/Sort/Quicksort.lean} \hfill quicksort y sus teoremas
\end{tcolorbox}
\end{minipage}

\navegacion{nivel1}{Nivel 1}{nivel3}{Nivel 3}

\newpage
%=====================================================================
% NIVEL 3
%=====================================================================
\nivel{3}{1}{\faCubes}{Datos: \texttt{inductive}, listas y conjuntos}{saber cómo se
escriben en Lean los objetos de la tesis: fórmulas, listas y contextos.}

\begin{minipage}[t]{0.335\textwidth}
\vspace{0pt}
{\large\bfseries\color{m1}\faCubes\ Fórmulas}\\[2pt]
\begin{lean}[]
inductive Formula : Type
| atom : String → Formula
| neg  : Formula → Formula
| impl : Formula → Formula → Formula
\end{lean}
Tres \textbf{constructores} $=$ las tres únicas formas de hacer una fórmula.\\[3pt]
\begin{center}
\begin{tikzpicture}[n/.style={draw=m1, fill=m1!8, rounded corners=4pt, font=\ttfamily\small, inner sep=4pt},
  level distance=1.25cm, sibling distance=3.4cm, edge from parent/.style={draw=m1, line width=1.1pt}]
  \node[n] {impl \ (⟶)}
    child {node[n] {neg \ (\textasciitilde)} child {node[n] {atom "P"}}}
    child {node[n] {atom "Q"}};
\end{tikzpicture}
\end{center}
\renewcommand{\arraystretch}{1.15}
\begin{tabular}{@{}l@{\ \ }l@{}}
papel: & $\neg P\to Q$\\
Lean: & \cod{~P ⟶ Q}\\
\end{tabular}\\[2pt]
sin notación:\\ {\footnotesize\cod{.impl (.neg (.atom "P")) (.atom "Q")}}
\end{minipage}\hfill
\begin{minipage}[t]{0.30\textwidth}
\vspace{0pt}
{\large\bfseries\color{m4b}\faListOl\ Listas} \ (\cod{List α})\\[4pt]
\renewcommand{\arraystretch}{1.35}
\begin{tabular}{@{}l@{\quad}l@{}}
\cod{[]} & la lista vacía\\
\cod{p :: cola} & $p$ delante de \cod{cola}\\
\cod{[3, 1, 2]} & \cod{3 :: 1 :: 2 :: []}\\
\cod{l₁ ++ l₂} & pegar dos listas\\
\cod{l.length} & longitud\\
\cod{l.filter q} & los que cumplen \cod{q}\\
\cod{a ∈ l} & $a$ aparece en \cod{l}\\
\end{tabular}

\vspace{4pt}
\clave{m4b}{Una lista también es inductiva: sus constructores son \cod{[]} y \cod{::}.
Por eso hay inducción con dos casos \chip{11}.}
\end{minipage}\hfill
\begin{minipage}[t]{0.335\textwidth}
\vspace{0pt}
{\large\bfseries\color{m4a}\faLayerGroup\ Contextos} \ (\cod{Set Formula})\\[4pt]
\renewcommand{\arraystretch}{1.35}
\begin{tabular}{@{}l@{\quad}l@{}}
\textbf{papel} & \textbf{Lean}\\
$\Gamma$ & \cod{Γ : Set Formula}\\
$\varnothing$ & \cod{(∅ : Set Formula)}\\
$\Gamma,\varphi$ & \cod{insert φ Γ}\\
$\varphi\in\Gamma$ & \cod{φ ∈ Γ}\\
$\Gamma\subseteq\Delta$ & \cod{Γ ⊆ Δ}\\
$\bot$ & \cod{Falsedad}\\
$\Gamma\vdash\varphi$ & \cod{ND Γ φ}\\
\end{tabular}

\vspace{4pt}
\clave{m4a}{\textbf{¿Por qué \cod{Set} y no \cod{List}?} En un contexto no importan
el orden ni las repeticiones.}
\end{minipage}

\pruebate{m1}{\item Escribe $\neg\neg P$ y $P\to(Q\to P)$ en Lean.
\item ¿Qué significa \cod{insert φ Γ}?
\item Escribe \cod{[1, 2]} solo con \cod{::} y \cod{[]}.}
{1.\ \cod{~~P} y \cod{P ⟶ (Q ⟶ P)}.
2.\ $\Gamma\cup\{\varphi\}$: el contexto $\Gamma$ con la hipótesis $\varphi$ añadida.
3.\ \cod{1 :: 2 :: []}.}

\navegacion{nivel2}{Nivel 2}{nivel4}{Nivel 4}

\newpage
%=====================================================================
% NIVEL 4
%=====================================================================
\nivel{4}{1}{\faCogs}{Funciones: \texttt{def}}{leer cómo se definen las funciones
de la tesis (evaluar, átomos, quicksort) y cómo se aplican.}

\begin{minipage}[t]{0.47\textwidth}
\vspace{0pt}
\begin{leanL}[title={Una función por casos: una línea por constructor}]
def evaluar (v : Valuacion) : Formula → Prop
| .atom s   => v s                       @1
| .neg p    => ¬ evaluar v p             @2
| .impl p q => evaluar v p → evaluar v q @3
\end{leanL}
\ex{1}{Un átomo vale lo que diga la valuación.}
\ex{2}{\cod{¬} es el «no» \textbf{de Lean}. La función se llama a sí misma.}
\ex{3}{\cod{→} es el «implica» \textbf{de Lean}.}
\begin{leanL}[]
def atomos : Formula → List String
| .atom s   => [s]
| .neg p    => atomos p
| .impl p q => atomos p ++ atomos q
\end{leanL}
\begin{leanL}[title={Thesis/Sort/Quicksort.lean}]
variable {α : Type _} [LinearOrder α]
def Ordenada (l : List α) : Prop :=
  l.Pairwise (· ≤ ·)
\end{leanL}
\end{minipage}\hfill
\begin{minipage}[t]{0.50\textwidth}
\vspace{0pt}
\begin{tcbraster}[raster columns=2, raster column skip=6pt, raster row skip=0pt, raster equal height=rows]
\tarj{m1}{Aplicar: sin paréntesis}{\cod{evaluar v φ} es $\mathrm{evaluar}(v,\varphi)$. Los paréntesis solo agrupan.}
\tarj{m1}{\cod{fun x => …}}{Función sin nombre ($\lambda$). \cod{(fun x => x + 1) 4} vale 5.}
\tarj{m1}{Notación punto}{\cod{l.length} es \cod{List.length l}. \cod{cola.filter q} es \cod{List.filter q cola}.}
\tarj{m1}{\cod{(· ≤ ·)}}{Abrevia \cod{fun a b => a ≤ b}: cada punto es un hueco.}
\tarj{m1}{\cod{_}}{«Dedúcelo tú»: Lean rellena lo que puede inferir.}
\tarj{m1}{\cod{noncomputable}}{Definición que no se puede ejecutar; \cod{literal} lo es porque decide una \cod{Prop}.}
\end{tcbraster}
\begin{tcolorbox}[colback=white, colframe=m1, boxrule=1.1pt, arc=5pt,
  title={\faSlidersH\ Tres clases de argumentos}, colbacktitle=m1, fonttitle=\bfseries,
  left=6pt, right=6pt, top=3pt, bottom=3pt]
\renewcommand{\arraystretch}{1.3}
\begin{tabular}{@{}l@{\quad}l@{\quad}l@{}}
\cod{(φ : Formula)} & \textbf{explícito} & lo escribes tú\\
\cod{{Γ φ}} & \textbf{implícito} & Lean lo deduce del resto\\
\cod{[LinearOrder α]} & \textbf{instancia} & Lean la busca sola\\
\end{tabular}
\end{tcolorbox}
\end{minipage}

\pruebate{m1}{\item ¿Qué devuelve \cod{atomos (P ⟶ P)}?
\item En \cod{evaluar}, ¿de quién son \cod{¬} y \cod{→}: de la fórmula o de Lean?
\item ¿Qué quiere decir \cod{{Γ φ}} entre llaves?}
{1.\ \cod{["P", "P"]}: con repetición (luego se usa \cod{dedup}).
2.\ De Lean: la fórmula usa \cod{~} y \cod{⟶}.
3.\ Argumentos implícitos: Lean los deduce.}

\navegacion{nivel3}{Nivel 3}{nivel5}{Nivel 5}

\newpage
%=====================================================================
% NIVEL 5
%=====================================================================
\nivel{5}{1}{\faQuoteRight}{Enunciar: \texttt{theorem}}{leer el enunciado de
cualquier teorema de la tesis y decirlo en español.}

\begin{center}
\begin{tikzpicture}[
  p/.style={draw=#1, fill=#1!10, line width=1.3pt, rounded corners=4pt, font=\ttfamily\normalsize,
    inner xsep=5pt, inner ysep=6pt},
  lab/.style={font=\normalsize, align=center, text width=#1}]
  \node[p=kw] (a) at (0,0) {theorem};
  \node[p=m1, right=1.5mm of a] (b) {completitud\_ND};
  \node[p=m4b, right=1.5mm of b] (c) {(φ : Formula) (h : EsTautologia φ)};
  \node[p=gris, right=1.5mm of c] (d) {:};
  \node[p=m4a, right=1.5mm of d] (e) {ND (∅ : Set Formula) φ};
  \node[p=tac, right=1.5mm of e] (f) {:= by};
  \node[lab=2.6cm, text=kw] at ($(a.south)+(0,-0.85)$) {palabra\\ clave};
  \node[lab=3.2cm, text=m1!80!black] at ($(b.south)+(0,-0.85)$) {nombre};
  \node[lab=8.6cm, text=m4b!80!black] at ($(c.south)+(0,-0.85)$) {\textbf{hipótesis}: «para toda fórmula $\varphi$,\\ si $h$ prueba que $\varphi$ es tautología…»};
  \node[lab=6.4cm, text=m4a!80!black] at ($(e.south)+(0,-0.85)$) {\textbf{enunciado}: «…entonces\\ $\varnothing\vdash\varphi$»};
  \node[lab=3.0cm, text=tac] at ($(f.south)+(0,-0.85)$) {empieza\\ la prueba};
\end{tikzpicture}
\end{center}

\begin{minipage}[t]{0.50\textwidth}
\vspace{0pt}
\begin{estado}[title={\#check @completitud\_ND}]
completitud_ND : ∀ (φ : Formula), EsTautologia φ → ND ∅ φ
\end{estado}
\clave{m1}{Lo que va \textbf{antes de los dos puntos} se convierte en $\forall$ y $\to$.
Un teorema es una \textbf{función}: recibe $\varphi$ y una prueba $h$, y devuelve una derivación.}
\begin{lean}[title={Usar un teorema = aplicarlo (Thesis/Prop/Main.lean)}]
theorem identidad_por_completitud :
    ND (∅ : Set Formula) tautologiaIdentidad :=
  completitud_ND tautologiaIdentidad identidad_es_tautologia
\end{lean}
\end{minipage}\hfill
\begin{minipage}[t]{0.47\textwidth}
\vspace{0pt}
\begin{tcbraster}[raster columns=2, raster column skip=6pt, raster row skip=0pt, raster equal height=rows]
\tarj{m1}{Modo término}{\cod{:= expresión}. La prueba es una sola aplicación de lemas.}
\tarj{tac}{Modo táctico}{\cod{:= by} y luego tácticas, paso a paso. Produce el mismo término.}
\tarj{m1}{\cod{theorem} y \cod{lemma}}{Son lo mismo para Lean. \cod{example} es un teorema sin nombre.}
\tarj{m1}{Nombres de la tesis}{\cod{h…}: un hecho. \cod{d…}: una derivación. \cod{hipInd}: hipótesis inductiva.}
\end{tcbraster}
\tarj{com}{Los comentarios de cada prueba}{\cod{-- A probar:} qué falta. \cod{-- Método:} cómo se hará. Léelos primero: son el resumen en español.}
\end{minipage}

\pruebate{m1}{\item Di en español: \cod{theorem debilitamiento {Γ φ} (d : ND Γ φ) : ∀ {Δ}, Γ ⊆ Δ → ND Δ φ}
\item ¿Qué devuelve \cod{completitud_ND φ h}?
\item ¿En qué se diferencian el modo término y el táctico?}
{1.\ Si $d$ deriva $\varphi$ desde $\Gamma$, entonces para todo $\Delta\supseteq\Gamma$ se deriva $\varphi$ desde $\Delta$.
2.\ Una derivación de $\varphi$ desde el contexto vacío.
3.\ Solo en cómo se escribe; el resultado es el mismo término.}

\navegacion{nivel4}{Nivel 4}{nivel6}{Nivel 6}

\newpage
%=====================================================================
% NIVEL 6
%=====================================================================
\nivel{6}{2}{\faEye}{El estado de prueba}{es lo que Lean te muestra en cada
paso. Leer el estado es leer la prueba.}

\begin{minipage}[t]{0.40\textwidth}
\vspace{0pt}
\begin{leanL}[title={Thesis/Prop/Main.lean}]
theorem identidad_es_tautologia :
    EsTautologia tautologiaIdentidad := by
  intro v                              @1
  change evaluar v P → evaluar v P     @2
  intro hP                             @3
  exact hP                             @4
\end{leanL}
\begin{estado}[title={Inicio: lo primero que muestra Lean}]
⊢ EsTautologia tautologiaIdentidad
\end{estado}
\end{minipage}\hfill
\begin{minipage}[t]{0.57\textwidth}
\vspace{0pt}
\begin{estadoL}[title={Anatomía de un estado}]
v : Valuacion        @@ arriba: lo que TIENES (hipótesis)
hP : evaluar v P     @@ cada una es  nombre : tipo
⊢ evaluar v P        @@ tras ⊢: lo que FALTA (la meta)
\end{estadoL}
\clave{m2}{Cada táctica \textbf{transforma el estado}. La prueba termina cuando no quedan metas.}
\end{minipage}

\vspace{2pt}
\vspace{6pt}
{\large\bfseries\color{m2}\faFilm\ La película: el estado después de cada paso}\\[1pt]
\begin{tcbraster}[raster columns=4, raster column skip=5pt, raster equal height]
\begin{estado}[title={Tras 1: intro v}]
v : Valuacion
⊢ evaluar v tautologiaIdentidad
\end{estado}
\begin{estado}[title={Tras 2: change …}]
v : Valuacion
⊢ evaluar v P → evaluar v P
\end{estado}
\begin{estado}[title={Tras 3: intro hP}]
v : Valuacion
hP : evaluar v P
⊢ evaluar v P
\end{estado}
\begin{tcolorbox}[estadobox, title={Tras 4: exact hP}, top=3pt, valign=center, halign=center]
{\color{verde}\Large\faCheckCircle}\ \ sin metas
\end{tcolorbox}
\end{tcbraster}
\begin{tcbraster}[raster columns=4, raster column skip=5pt, raster equal height,
  raster every box/.style={colback=m2!6, colframe=m2!50, boxrule=0.6pt, arc=3pt,
    left=4pt, right=4pt, top=2pt, bottom=2pt, before upper=\raggedright\small}]
\begin{tcolorbox}\bola{1} \cod{EsTautologia} es un $\forall v$: \cod{intro} toma una $v$ cualquiera.\end{tcolorbox}
\begin{tcolorbox}\bola{2} La misma meta, desplegada por definición. Ya se ve el $\to$.\end{tcolorbox}
\begin{tcolorbox}\bola{3} La meta es un $\to$: \cod{intro} supone el antecedente.\end{tcolorbox}
\begin{tcolorbox}\bola{4} \cod{hP} es justo la meta: \cod{exact} la cierra.\end{tcolorbox}
\end{tcbraster}

\dice{«Sea $v$ una valuación. Por definición hay que probar que si $P$ vale en $v$, entonces $P$ vale en $v$. Supongámoslo: es justo lo que había que probar.»}

\pruebate{m2}{\item En un estado, ¿qué hay arriba y qué hay tras \cod{⊢}?
\item ¿Cuándo termina una prueba?
\item ¿Por qué el primer \cod{intro} introduce una valuación?}
{1.\ Arriba, las hipótesis disponibles; tras \cod{⊢}, la meta.
2.\ Cuando no quedan metas.
3.\ Porque \cod{EsTautologia φ} es por definición \cod{∀ v, evaluar v φ}.}

\navegacion{nivel5}{Nivel 5}{nivel7}{Nivel 7}

\newpage
%=====================================================================
% NIVEL 7
%=====================================================================
\nivel{7}{2}{\faHandPointRight}{Cinco tácticas que mueven todo}{con
\texttt{intro}, \texttt{exact}, \texttt{apply}, \texttt{have} y \texttt{constructor} se
escribe casi toda la tesis.}

% #1 táctica, #2 qué hace, #3 se dice
\newcommand{\tact}[3]{\begin{tcolorbox}[colback=m2!6, colframe=m2, boxrule=1pt, arc=4pt,
  left=5pt, right=5pt, top=2pt, bottom=2pt, leftrule=5pt, before upper=\raggedright,
  before skip=0pt, after skip=0pt, height=2.42cm, valign=center]
  {\large\ttfamily\bfseries\color{tac}#1}\quad #2\\[1pt]
  {\small\color{gris}\faCommentDots\ #3}\end{tcolorbox}}

\begin{minipage}[t]{0.43\textwidth}\vspace{0pt}
\tact{intro hA}{Saca el antecedente de la meta y lo pone arriba como hipótesis.}{«Supongamos $A$.» / «Sea $x$ cualquiera.»}
\end{minipage}\hfill
\begin{minipage}[t]{0.27\textwidth}\vspace{0pt}
\begin{estadoS}[title={Antes}, height=2.42cm, before skip=0pt, after skip=0pt]
⊢ A → B
\end{estadoS}
\end{minipage}\hfill
\begin{minipage}[t]{0.27\textwidth}\vspace{0pt}
\begin{estadoS}[title={Después}, height=2.42cm, before skip=0pt, after skip=0pt]
hA : A
⊢ B
\end{estadoS}
\end{minipage}

\vspace{3pt}
\begin{minipage}[t]{0.43\textwidth}\vspace{0pt}
\tact{exact e}{Cierra la meta: \cod{e} tiene exactamente ese tipo.}{«…que es justo lo que había que probar.»}
\end{minipage}\hfill
\begin{minipage}[t]{0.27\textwidth}\vspace{0pt}
\begin{estadoS}[title={Antes}, height=2.42cm, before skip=0pt, after skip=0pt]
hA : A
⊢ A
\end{estadoS}
\end{minipage}\hfill
\begin{minipage}[t]{0.27\textwidth}\vspace{0pt}
\begin{tcolorbox}[estadobox, title={Después de \texttt{exact hA}}, height=2.42cm, before skip=0pt,
  after skip=0pt, top=3pt, valign=center, halign=center]
{\color{verde}\Large\faCheckCircle}\ \ sin metas
\end{tcolorbox}
\end{minipage}

\vspace{3pt}
\begin{minipage}[t]{0.43\textwidth}\vspace{0pt}
\tact{apply h}{Razona hacia atrás: si \cod{h : A → B} y la meta es \cod{B}, ahora basta \cod{A}.}{«Por $h$, basta probar $A$.»}
\end{minipage}\hfill
\begin{minipage}[t]{0.27\textwidth}\vspace{0pt}
\begin{estadoS}[title={Antes}, height=2.42cm, before skip=0pt, after skip=0pt]
hAB : A → B
⊢ B
\end{estadoS}
\end{minipage}\hfill
\begin{minipage}[t]{0.27\textwidth}\vspace{0pt}
\begin{estadoS}[title={Después de apply hAB}, height=2.42cm, before skip=0pt, after skip=0pt]
hAB : A → B
⊢ A
\end{estadoS}
\end{minipage}

\vspace{3pt}
\begin{minipage}[t]{0.43\textwidth}\vspace{0pt}
\tact{have hB : B := e}{Razona hacia adelante: añade un hecho nuevo con nombre. La meta no cambia.}{«Primero observamos que $B$.»}
\end{minipage}\hfill
\begin{minipage}[t]{0.27\textwidth}\vspace{0pt}
\begin{estadoS}[title={Antes}, height=2.42cm, before skip=0pt, after skip=0pt]
hAB : A → B
hA : A
⊢ B
\end{estadoS}
\end{minipage}\hfill
\begin{minipage}[t]{0.27\textwidth}\vspace{0pt}
\begin{estadoS}[title={Tras have hB : B := hAB hA}, height=2.42cm, before skip=0pt, after skip=0pt]
hAB : A → B
hA : A
hB : B
⊢ B
\end{estadoS}
\end{minipage}

\vspace{3pt}
\begin{minipage}[t]{0.43\textwidth}\vspace{0pt}
\tact{constructor}{Parte la meta en sus componentes. \cod{refine ⟨?_, ?_⟩} hace lo mismo: cada \cod{?_} es una meta nueva.}{«Probamos cada parte por separado.»}
\end{minipage}\hfill
\begin{minipage}[t]{0.27\textwidth}\vspace{0pt}
\begin{estadoS}[title={Antes}, height=2.42cm, before skip=0pt, after skip=0pt]
⊢ A ∧ B
\end{estadoS}
\end{minipage}\hfill
\begin{minipage}[t]{0.27\textwidth}\vspace{0pt}
\begin{estadoS}[title={Después: dos metas}, height=2.42cm, before skip=0pt, after skip=0pt]
case left
⊢ A
case right
⊢ B
\end{estadoS}
\end{minipage}

\vspace{5pt}
\begin{tcbraster}[raster columns=2, raster column skip=6pt, raster equal height]
\clave{m2}{\textbf{El punto \cod{·}} enfoca una sola meta. Tras \cod{constructor} verás dos \cod{·}, uno por parte. {\small\color{gris}(Estados comprobados en Lean; se omite la línea \cod{A B : Prop}.)}}
\clave{m2}{\textbf{Hacia atrás o hacia adelante:} \cod{apply} cambia la meta; \cod{have} añade hipótesis. La tesis usa casi siempre \cod{have}.}
\end{tcbraster}
\navegacion{nivel6}{Nivel 6}{nivel8}{Nivel 8}

\newpage
%=====================================================================
% NIVEL 8
%=====================================================================
\nivel{8}{2}{\faTable}{La tabla maestra: probar y usar}{ante cualquier símbolo
lógico, saber qué táctica toca. Mira si está en la \textbf{meta} o en una \textbf{hipótesis}.}

\begin{center}
\renewcommand{\arraystretch}{1.34}
\begin{tabular}{@{}>{\large}l@{\quad}>{\raggedright}p{7.3cm}@{\quad}>{\raggedright}p{8.0cm}@{\quad}>{\raggedright\arraybackslash\leavevmode\color{gris}}p{6.6cm}@{}}
 & \cellcolor{m2!18}\textbf{\faBullseye\ Está en la meta: para \emph{probarlo}} & \cellcolor{m4b!18}\textbf{\faToolbox\ Es una hipótesis \texttt{h}: para \emph{usarlo}} & \textbf{\faCommentDots\ Se dice}\\
\cod{A → B} & \cod{intro hA} & \cod{h hA} \ o \ \cod{apply h} & «supongamos $A$» / «por $h$»\\
\cod{∀ x, P x} & \cod{intro x} & \cod{h a} & «sea $x$ cualquiera» / «en particular»\\
\cod{A ∧ B} & \cod{constructor} \ o \ \cod{⟨hA, hB⟩} & \cod{obtain ⟨hA, hB⟩ := h} \ o \ \cod{h.left}, \cod{h.right} & «las dos partes»\\
\cod{A ∨ B} & \cod{Or.inl hA} \ o \ \cod{Or.inr hB} & \cod{rcases h with hA | hB} & «por casos»\\
\cod{A ↔ B} & \cod{constructor} & \cod{h.mp} (ida), \cod{h.mpr} (vuelta) & «ida y vuelta»\\
\cod{¬ A} & \cod{intro hA} \ (queda probar \cod{False}) & \cod{h hA} \ (da \cod{False}) & «supongamos $A$: contradicción»\\
\cod{∃ x, P x} & \cod{⟨a, ha⟩} & \cod{obtain ⟨x, hx⟩ := h} & «tomamos un testigo»\\
\cod{False} & — & \cod{False.elim h} & «de lo falso sale cualquier cosa»\\
\cod{a = b} & \cod{rfl} & \cod{rw [h]} \ o \ \cod{subst} \chip{9} & «reescribiendo»\\
\end{tabular}
\end{center}

\begin{tcbraster}[raster columns=3, raster column skip=6pt, raster equal height]
\clave{m2}{\textbf{\cod{¬ A} es por definición \cod{A → False}.} Por eso se prueba con \cod{intro} y se usa aplicándola.}
\clave{m2}{\textbf{\cod{⟨ , ⟩}} es el constructor anónimo: arma un par (para \cod{∧}, \cod{∃}, \cod{↔}) sin decir su nombre.}
\clave{m2}{Es Curry–Howard en acción: usar \cod{A → B} es \textbf{aplicar una función}; usar \cod{A ∧ B} es \textbf{abrir un par}.}
\end{tcbraster}

\pruebate{m2}{\item La meta es \cod{evaluar v φ → ND Γ φ}. ¿Qué táctica?
\item Tienes \cod{h : A ↔ B} y \cod{hA : A}. ¿Cómo obtienes \cod{B}?
\item Tienes \cod{hNeg : ¬ A} y \cod{hA : A}. ¿Qué es \cod{hNeg hA}?}
{1.\ \cod{intro hev}.
2.\ \cod{h.mp hA}.
3.\ Una prueba de \cod{False}.}

\navegacion{nivel7}{Nivel 7}{nivel9}{Nivel 9}

\newpage
%=====================================================================
% NIVEL 9
%=====================================================================
\nivel{9}{2}{\faEquals}{Reescribir y calcular}{cambiar la forma de la meta sin
cambiar lo que dice: desplegar definiciones, sustituir iguales, encadenar pasos.}

\begin{minipage}[t]{0.42\textwidth}
\vspace{0pt}
\tarj{tac}{\cod{rw [h]}}{Si \cod{h : a = b}, cambia \cod{a} por \cod{b} en la meta. Con \cod{at h2}, en la hipótesis \cod{h2}. También despliega definiciones.}
\tarj{tac}{\cod{simp only [lemas]} \ y \ \cod{simpa … using h}}{Reescribe con esa lista de lemas hasta que no puede más. \cod{simpa} además cierra la meta con \cod{h}.}
\tarj{tac}{\cod{change P}}{Cambia la meta por \cod{P}, que es \textbf{la misma por definición}. No prueba nada: la muestra desplegada.}
\tarj{tac}{\cod{subst x} \ y \ \cod{rfl}}{\cod{subst}: si \cod{x = a}, sustituye \cod{x} en todas partes. \cod{rfl}: los dos lados son iguales al calcular.}
\tarj{tac}{\cod{calc}}{Cadena de pasos unidos por transitividad. Cada \cod{_} es «lo anterior».}
\end{minipage}\hfill
\begin{minipage}[t]{0.55\textwidth}
\vspace{0pt}
\begin{tcbraster}[raster columns=2, raster column skip=5pt, raster equal height]
\begin{estado}[title={Antes de rw [h1]}]
h1 : a = b
h2 : b = c
⊢ a = c
\end{estado}
\begin{estado}[title={Después}]
h1 : a = b
h2 : b = c
⊢ b = c
\end{estado}
\end{tcbraster}
\begin{lean}[title={Reescribir en una hipótesis (kalmar, caso atómico)}]
rw [literal, if_pos hvs] at hmem   -- despliega literal
\end{lean}
\begin{estado}[title={longitud\_filtro\_le: la meta antes y después de simp only […]}]
⊢ (filter q (a :: cola)).length ≤ (a :: cola).length
⊢ (filter q cola).length + 1 ≤ cola.length + 1
\end{estado}
\begin{lean}[title={calc: la cadena de quicksort\_permutacion}]
calc
  quicksort menoresIguales ++ (p :: quicksort mayores)
      ~ menoresIguales ++ (p :: mayores) := hConcatenacion
  _ ~ p :: (menoresIguales ++ mayores) := hPivote
  _ ~ p :: cola := hCons
\end{lean}
\end{minipage}

\pruebate{m2}{\item ¿Qué diferencia hay entre \cod{rw} y \cod{change}?
\item ¿Qué significa \cod{at hmem}?
\item En un \cod{calc}, ¿qué es \cod{_}?}
{1.\ \cod{rw} usa una igualdad o ecuación; \cod{change} solo despliega: la meta es la misma por definición.
2.\ Que se reescribe en la hipótesis \cod{hmem}, no en la meta.
3.\ El lado derecho del paso anterior.}

\navegacion{nivel8}{Nivel 8}{nivel10}{Nivel 10}

\newpage
%=====================================================================
% NIVEL 10
%=====================================================================
\nivel{10}{2}{\faCodeBranch}{Casos y absurdo. Nivel objeto y nivel meta}{separar
la lógica \textbf{de Lean} (con la que razonas) de la lógica \textbf{del cálculo} (sobre la que razonas).}

\begin{minipage}[t]{0.40\textwidth}
\vspace{0pt}
\tarj{tac}{\cod{by_cases hp : P}}{Dos metas: una con \cod{hp : P} y otra con \cod{hp : ¬P}.}
\begin{estado}[title={Tras by\_cases hA : A}]
case pos
hA : A
⊢ B
case neg
hA : ¬A
⊢ B
\end{estado}
\tarj{tac}{\cod{rcases h with hA | hB}}{Si \cod{h : A ∨ B}: un caso por lado. \cod{cases hq : q a with} hace casos sobre un dato y \cod{hq} recuerda cuál.}
\tarj{tac}{\cod{by_contra h}}{Supón la negación de la meta; ahora hay que probar \cod{False}.}
\end{minipage}\hfill
\begin{minipage}[t]{0.57\textwidth}
\vspace{0pt}
\begin{tcolorbox}[colback=white, colframe=black!75, boxrule=1.3pt, arc=6pt,
  title={\faBalanceScale\ La distinción que más preguntan}, colbacktitle=black!75,
  fonttitle=\bfseries\large, left=4pt, right=4pt, top=2pt, bottom=2pt]
\renewcommand{\arraystretch}{1.32}
\begin{tabular}{@{}>{\bfseries}l@{\ \ }>{\raggedright}p{5.6cm}@{\ \ }>{\raggedright\arraybackslash}p{5.3cm}@{}}
 & \cellcolor{m2!18}\textbf{META: la lógica de Lean} & \cellcolor{m4a!18}\textbf{OBJETO: el cálculo ND}\\
negación & \cod{¬ P} & \cod{~φ}\\
implicación & \cod{P → Q} & \cod{φ ⟶ ψ}\\
lo falso & \cod{False} & \cod{Falsedad}\\
demostrable & hay un término de tipo \cod{P} & \cod{ND Γ φ}\\
absurdo & \cod{by_contra} & \cod{ND.classical} (regla RAA)\\
casos & \cod{by_cases} & \cod{analisis_por_casos}\\
\end{tabular}
\end{tcolorbox}
\begin{lean}[title={Los dos niveles, en dos líneas reales de la tesis}]
by_contra hnp                  -- META: absurdo de Lean (correccion)
exact ND.classical dFalsedad   -- OBJETO: regla RAA (debilitamiento)
\end{lean}
\clave{m2}{\textbf{Consecuencia:} \cod{by_cases} y \cod{by_contra} sobre proposiciones cualesquiera
son lógica clásica de Lean. Por eso la completitud depende de \cod{Classical.choice} \chip{19}.}
\end{minipage}

\pruebate{m2}{\item ¿\cod{ND.classical} usa lógica clásica de Lean?
\item ¿Qué tienes después de \cod{by_contra h}?
\item ¿Cuál es la versión objeto de \cod{False}?}
{1.\ No: es un constructor, un dato. La regla es clásica \emph{dentro} del cálculo.
2.\ \cod{h}: la negación de la meta; y la meta nueva es \cod{False}.
3.\ \cod{Falsedad}, la fórmula \cod{~(⊥ ⟶ ⊥)} con un átomo llamado «⊥».}

\navegacion{nivel9}{Nivel 9}{nivel11}{Nivel 11}

\newpage
%=====================================================================
% NIVEL 11
%=====================================================================
\nivel{11}{3}{\faListOl}{Inducción sobre datos}{probar algo para todas las listas
o todas las fórmulas: un caso por constructor.}

\begin{minipage}[t]{0.52\textwidth}
\vspace{0pt}
\begin{lean}[title={Thesis/Sort/Quicksort.lean (el caso false, recortado)}]
theorem longitud_filtro_le (q : α → Bool) (l : List α) :
    (l.filter q).length ≤ l.length := by
  induction l with                                  @1
  | nil => exact Nat.le_refl 0                      @2
  | cons a cola hipInd =>                           @3
      cases hq : q a with                           @4
      | true =>
          simp only [List.filter_cons, hq, ↓reduceIte,
            List.length_cons]
          exact Nat.succ_le_succ hipInd             @5
      | false => …
\end{lean}
\begin{estado}[title={Estado real al entrar en el caso cons, rama true}]
a : α
cola : List α
hipInd : (filter q cola).length ≤ cola.length
hq : q a = true
⊢ (filter q (a :: cola)).length ≤ (a :: cola).length
\end{estado}
\end{minipage}\hfill
\begin{minipage}[t]{0.45\textwidth}
\vspace{0pt}
\ex{1}{\cod{induction l with}: un caso por constructor de la lista.}
\ex{2}{\cod{nil}: la lista vacía. Caso base.}
\ex{3}{\cod{cons a cola hipInd}: tú pones los nombres de la cabeza, la cola y la \textbf{hipótesis inductiva}.}
\ex{4}{Dentro, casos sobre si \cod{a} pasa el filtro.}
\ex{5}{Aquí se usa \cod{hipInd}: sumar 1 a los dos lados.}
\clave{m3}{\textbf{\cod{hipInd} es el mismo enunciado, pero para \cod{cola}.} Mira el estado:
habla de \cod{cola}; la meta habla de \cod{a :: cola}.}
\begin{leanS}[title={Sobre fórmulas (kalmar): tres constructores, tres casos}]
induction φ with
| atom s => …                    -- sin hipótesis inductiva
| neg p hipInd => …              -- una: para p
| impl p q hipIndP hipIndQ => …  -- dos: para p y para q
\end{leanS}
\end{minipage}

\pruebate{m3}{\item ¿Cuántas hipótesis inductivas tiene el caso \cod{impl}? ¿Por qué?
\item En el caso \cod{cons a cola hipInd}, ¿qué dice \cod{hipInd}?
\item ¿Para qué sirve el \cod{hq} de \cod{cases hq : q a}?}
{1.\ Dos: \cod{impl} tiene dos subfórmulas.
2.\ El enunciado del teorema, para \cod{cola}.
3.\ Guarda la ecuación \cod{q a = true} o \cod{q a = false} del caso.}

\navegacion{nivel10}{Nivel 10}{nivel12}{Nivel 12}

\newpage
%=====================================================================
% NIVEL 12
%=====================================================================
\nivel{12}{3}{\faProjectDiagram}{Inducción sobre reglas (derivaciones)}{es la
técnica central del capítulo 11: las derivaciones también son datos inductivos.}

\begin{minipage}[t]{0.50\textwidth}
\vspace{0pt}
\begin{lean}[title={Thesis/Prop/NaturalDeduction.lean: el cálculo completo}]
inductive ND : Set Formula → Formula → Prop
| hyp {Γ φ} : φ ∈ Γ → ND Γ φ                              @1
| impI {Γ φ ψ} : ND (insert φ Γ) ψ → ND Γ (φ ⟶ ψ)        @2
| impE {Γ φ ψ} : ND Γ (φ ⟶ ψ) → ND Γ φ → ND Γ ψ          @3
| negI {Γ φ} : ND (insert φ Γ) Falsedad → ND Γ (~φ)       @4
| negE {Γ φ} : ND Γ (~φ) → ND Γ φ → ND Γ Falsedad         @5
| botE {Γ φ} : ND Γ Falsedad → ND Γ φ                     @6
| classical {Γ φ} : ND (insert (~φ) Γ) Falsedad → ND Γ φ  @7
\end{lean}
\begin{tcolorbox}[colback=m3!6, colframe=m3, boxrule=1pt, arc=4pt, left=6pt, right=6pt,
  top=3pt, bottom=3pt]
\bola{1} hip. \ \bola{2} $\to$I \ \bola{3} $\to$E \ \bola{4} $\neg$I \ \bola{5} $\neg$E \
\bola{6} $\bot$E \ \bola{7} RAA\\[3pt]
\textbf{Cómo se lee un constructor:} lo que hay antes de la última \cod{→} son las
\textbf{premisas}; lo de después, la \textbf{conclusión}. \cod{insert φ Γ} es la hipótesis que
la regla descarga.
\end{tcolorbox}
\clave{m3}{\textbf{7 constructores $=$ 7 reglas $=$ 7 casos} en toda inducción sobre una derivación.}
\ex{8}{El enunciado deja \cod{∀ {Δ}} \textbf{después} de los dos puntos. Así \cod{hipInd} vale para \emph{cualquier} $\Delta$ (generalizar antes de inducir).}
\ex{9}{El \cod{@} deja nombrar también los implícitos. \cod{_} ignora la subderivación; \cod{hipInd} es lo que se usa.}
\end{minipage}\hfill
\begin{minipage}[t]{0.47\textwidth}
\vspace{0pt}
\begin{leanS}[title={debilitamiento: el caso de la regla →I}]
theorem debilitamiento {Γ φ} (d : ND Γ φ) :
    ∀ {Δ}, Γ ⊆ Δ → ND Δ φ := by                      @8
  induction d with
  | @impI Γ φ ψ _ hipInd =>                          @9
      intro Δ hsub
      have hAmpliado : insert φ Γ ⊆ insert φ Δ :=
        Set.insert_subset_insert hsub
      have dCuerpo : ND (insert φ Δ) ψ := hipInd hAmpliado
      exact ND.impI dCuerpo
  | …   -- y los otros seis casos
\end{leanS}
\begin{estado}[title={Estado real tras «intro Δ hsub» (sin las variables con ✝)}]
hipInd : ∀ {Δ : Set Formula}, insert φ Γ ⊆ Δ → ND Δ ψ
Δ : Set Formula
hsub : Γ ⊆ Δ
⊢ ND Δ (φ ⟶ ψ)
\end{estado}
\dice{«Por inducción sobre la derivación. Caso $\to$I: la premisa deriva $\psi$ desde
$\Gamma,\varphi$. Como $\Gamma,\varphi\subseteq\Delta,\varphi$, la hipótesis inductiva da $\psi$
desde $\Delta,\varphi$, y $\to$I concluye.»}
\end{minipage}

\navegacion{nivel11}{Nivel 11}{nivel13}{Nivel 13}

\newpage
%=====================================================================
% NIVEL 13
%=====================================================================
\nivel{13}{3}{\faSync*}{Recursión bien fundada}{entender por qué quicksort
necesita \texttt{termination\_by} y cómo un teorema «se llama a sí mismo».}

\begin{minipage}[t]{0.58\textwidth}
\vspace{0pt}
\begin{leanS}[title={Thesis/Sort/Quicksort.lean (las pruebas de decreasing\_by, recortadas)}]
def quicksort : List α → List α
  | [] => []
  | p :: cola =>
      quicksort (cola.filter (fun x => decide (x ≤ p)))            @1
        ++ p :: quicksort (cola.filter (fun x => ! decide (x ≤ p)))
  termination_by l => l.length                                     @2
  decreasing_by                                                    @3
    · … exact longitud_filtro_lt_cons (fun x => decide (x ≤ p)) p cola
    · … exact longitud_filtro_lt_cons (fun x => ! decide (x ≤ p)) p cola
\end{leanS}
\begin{estadoS}[title={Una de las dos obligaciones (estado real)}]
⊢ (filter (fun x => decide (x ≤ p)) cola).length < (p :: cola).length
\end{estadoS}
\begin{leanS}[title={En un teorema, la llamada recursiva ES la hipótesis inductiva}]
theorem quicksort_permutacion : ∀ l : List α, quicksort l ~ l
  | [] => by …
  | p :: cola => by
      …
      have hipIndMayores : quicksort mayores ~ mayores :=
        quicksort_permutacion mayores                              @4
      …
  termination_by l => l.length
  decreasing_by
    · exact hMenoresIgualesLt
    · exact hMayoresLt
\end{leanS}
\end{minipage}\hfill
\begin{minipage}[t]{0.39\textwidth}
\vspace{0pt}
\ex{1}{La llamada no es sobre \cod{cola}, sino sobre un \textbf{filtro} de \cod{cola}: Lean no ve solo que sea más pequeño.}
\ex{2}{\cod{termination_by}: la \textbf{medida} que baja. Aquí, la longitud.}
\ex{3}{\cod{decreasing_by}: una obligación por cada llamada recursiva: probar que la medida baja.}
\ex{4}{El teorema se aplica a una lista más corta. Las mismas dos cotas lo justifican.}
\tarj{m3}{\cod{decide (x ≤ p)}}{Convierte la proposición \cod{x ≤ p} en un \cod{Bool}, que es lo que \cod{filter} necesita. El \cod{!} es el «no» booleano.}
\begin{tcolorbox}[colback=m3!6, colframe=m3, boxrule=1pt, arc=4pt, left=5pt, right=5pt,
  top=3pt, bottom=3pt, title={\faBalanceScale\ Dos recursiones}, colbacktitle=m3, fonttitle=\bfseries]
\textbf{Estructural:} sobre un trozo del dato (\cod{cola}). Lean la acepta sola: \cod{evaluar}, \cod{atomos}.\\[2pt]
\textbf{Bien fundada:} sobre algo más pequeño según una medida. Hay que probarlo: \cod{quicksort}.
\end{tcolorbox}
\begin{lean}[]
#eval quicksort [3, 1, 2]  -- [1, 2, 3]
\end{lean}
\end{minipage}

\pruebate{m3}{\item ¿Por qué Lean no acepta quicksort sin \cod{termination_by}?
\item ¿Cuántas obligaciones hay en \cod{decreasing_by}? ¿Por qué?
\item ¿Dónde está la hipótesis inductiva de \cod{quicksort_permutacion}?}
{1.\ La recursión no es estructural: se llama sobre un filtro, no sobre la cola.
2.\ Dos: una por cada llamada recursiva.
3.\ Es la llamada al propio teorema sobre \cod{menoresIguales} y \cod{mayores}.}

\navegacion{nivel12}{Nivel 12}{nivel14}{Nivel 14}

\newpage
%=====================================================================
% NIVEL 14
%=====================================================================
\nivel{14}{3}{\faTree}{Construir derivaciones}{ver que un árbol de deducción
natural y una prueba en Lean son el mismo objeto, línea por línea.}

\begin{minipage}[t]{0.40\textwidth}
\vspace{0pt}
{\large\bfseries\color{m3}El árbol de $\vdash\neg\neg P\to P$}\\[2pt]
{\small con $\Gamma=\{\neg\neg P\}$ y $\Delta=\Gamma\cup\{\neg P\}$}\\[4pt]
\begin{center}
\begin{tikzpicture}[
  n/.style={draw=m3, fill=m3!7, line width=1.1pt, rounded corners=4pt, inner sep=5pt, font=\normalsize},
  up/.style={-{Stealth[length=2.8mm]}, line width=1.3pt, m3}]
  \node[n] (a) at (0,0) {\bola{1}\ $\Delta\vdash\neg\neg P$};
  \node[n] (b) at (5.0,0) {\bola{2}\ $\Delta\vdash\neg P$};
  \node[n] (c) at (2.5,-1.75) {\bola{3}\ $\Delta\vdash\bot$};
  \node[n] (d) at (2.5,-3.5) {\bola{4}\ $\Gamma\vdash P$};
  \node[n, fill=m3!20] (e) at (2.5,-5.25) {\bola{5}\ $\varnothing\vdash\neg\neg P\to P$};
  \draw[up] (a) -- (c); \draw[up] (b) -- (c);
  \draw[up] (c) -- node[right, font=\small, align=left]{RAA\\ descarga $\neg P$} (d);
  \draw[up] (d) -- node[right, font=\small, align=left]{$\to$I\\ descarga $\neg\neg P$} (e);
  \node[font=\small, text=m3!70!black] at (0,0.75) {hipótesis};
  \node[font=\small, text=m3!70!black] at (5.0,0.75) {hipótesis};
  \node[font=\small, text=m3!70!black, anchor=east] at (2.0,-0.95) {$\neg$E};
\end{tikzpicture}
\end{center}
\clave{m3}{\textbf{Receta de la tesis:} \cod{let} nombra los contextos; los \cod{h…} prueban
pertenencias; cada \cod{d…} es \textbf{un nodo del árbol}; \cod{exact} pone la raíz.}
\end{minipage}\hfill
\begin{minipage}[t]{0.57\textwidth}
\vspace{0pt}
\begin{lean}[title={Thesis/Prop/Main.lean}]
theorem ej_eliminacion_doble_negacion :
    ND (∅ : Set Formula) ((~~P) ⟶ P) := by
  let Γ := insert (~~P) (∅ : Set Formula)
  let Δ := insert (~P) Γ
  have hDoble : (~~P) ∈ Γ := Set.mem_insert _ _
  have hDobleDelta : (~~P) ∈ Δ := Set.mem_insert_of_mem _ hDoble
  have hNeg : (~P) ∈ Δ := Set.mem_insert _ _
  have dDoble : ND Δ (~~P) := ND.hyp hDobleDelta           @1
  have dNeg : ND Δ (~P) := ND.hyp hNeg                     @2
  have dFalsedad : ND Δ Falsedad := ND.negE dDoble dNeg    @3
  have dP : ND Γ P := ND.classical dFalsedad               @4
  exact ND.impI dP                                         @5
\end{lean}
\begin{tcbraster}[raster columns=2, raster column skip=6pt, raster equal height]
\clave{m3}{El árbol se lee de las hojas a la raíz; Lean, de arriba abajo. \textbf{Es el mismo orden.}}
\clave{m3}{\cod{ND.negE dDoble dNeg}: una regla es una función que recibe las derivaciones de arriba y devuelve la de abajo.}
\end{tcbraster}
\begin{estadoS}[title={\#print axioms ej\_eliminacion\_doble\_negacion}]
'Thesis.Prop.ej_eliminacion_doble_negacion' does not depend on any axioms
\end{estadoS}
\clave{m3}{\textbf{Ningún axioma}, aunque usa RAA: \cod{ND.classical} es un constructor (un dato),
no lógica clásica de Lean.}
\end{minipage}

\navegacion{nivel13}{Nivel 13}{nivel15}{Nivel 15}

\newpage
%=====================================================================
% NIVEL 15
%=====================================================================
\nivel{15}{4}{\faPuzzlePiece}{Completitud: las piezas}{antes de leer el teorema,
conocer sus dos lemas grandes y saber leer sus enunciados.}

\begin{center}
\begin{tikzpicture}[
  b/.style={draw=#1, fill=#1!9, line width=1.2pt, rounded corners=5pt, font=\ttfamily\small,
    inner sep=5pt, align=center},
  f/.style={-{Stealth[length=2.8mm]}, line width=1.4pt, black!50}]
  \node[b=m1] (at) at (0,0) {atomos};
  \node[b=m1] (li) at (3.3,0) {literal};
  \node[b=m1] (cl) at (7.6,0) {contextoLiterales};
  \node[b=m4a] (pl) at (12.9,0) {pertenencia\_literal};
  \node[b=m4a, line width=2pt] (ka) at (17.6,0) {\textbf{kalmar}};
  \node[b=m3] (de) at (3.3,-1.5) {debilitamiento};
  \node[b=m3] (dn) at (9.2,-1.5) {introduccion\_doble\_negacion};
  \node[b=m3] (ac) at (15.6,-1.5) {analisis\_por\_casos};
  \node[b=m4a] (co) at (3.3,-3.0) {congruencia\_contextoLiterales};
  \node[b=m4a, line width=2pt] (ds) at (11.4,-3.0) {\textbf{descarga}};
  \node[b=m4a, fill=m4a!30, line width=2pt] (ct) at (20.0,-3.0) {\textbf{completitud\_ND}};
  \draw[f] (li) -- (cl); \draw[f] (cl) -- (pl); \draw[f] (pl) -- (ka);
  \draw[f] (de) -- (dn); \draw[f] (dn) -- (ka);
  \draw[f] (co) -- (ds); \draw[f] (ac) -- (ds);
  \draw[f] (ka) -| (ct); \draw[f] (ds) -- (ct);
  \draw[f] (at.north) -- ++(0,0.4) -| (ka.north);
  \node[font=\small, text=m1, anchor=west] at (-1.0,-0.72) {definiciones};
  \node[font=\small, text=m3, anchor=east] at (2.0,-2.2) {reglas derivadas (NaturalDeduction.lean)};
\end{tikzpicture}
\end{center}

\begin{minipage}[t]{0.50\textwidth}
\vspace{0pt}
\begin{lean}[title={El enunciado de Kalmár (Lema 11.15), con saltos de línea añadidos}]
theorem kalmar (v : Valuacion) (listaAtomos : List String) :
    ∀ φ : Formula, atomos φ ⊆ listaAtomos →               @1
      (evaluar v φ →
          ND (contextoLiterales v listaAtomos) φ) ∧       @2
      (¬ evaluar v φ →
          ND (contextoLiterales v listaAtomos) (~φ))      @3
\end{lean}
\ex{1}{«Fijada una fila $v$ y una lista que \textbf{cubre} los átomos de $\varphi$…»}
\ex{2}{«…si $\varphi$ es verdadera en $v$, los literales derivan $\varphi$…»}
\ex{3}{«…y si es falsa, derivan $\neg\varphi$.» Las dos mitades van unidas por \cod{∧}.}
\end{minipage}\hfill
\begin{minipage}[t]{0.47\textwidth}
\vspace{0pt}
\begin{lean}[title={El enunciado de la descarga (Lema 11.17)}]
theorem descarga {φ : Formula} :
    ∀ (listaAtomos : List String), listaAtomos.Nodup →  @4
      (∀ v, ND (contextoLiterales v listaAtomos) φ) →   @5
      ND (∅ : Set Formula) φ                            @6
\end{lean}
\ex{4}{«Si la lista no tiene repetidos…»}
\ex{5}{«…y \textbf{toda} fila deriva $\varphi$ desde sus literales…»}
\ex{6}{«…entonces $\varphi$ se deriva sin hipótesis.»}
\end{minipage}

\vspace{2pt}
\begin{tcolorbox}[colback=m4a!6, colframe=m4a, boxrule=1pt, arc=4pt, left=6pt, right=6pt,
  top=3pt, bottom=3pt, title={\faExchange*\ Del papel a Lean: número del resultado y nombre}, colbacktitle=m4a,
  fonttitle=\bfseries]
\small
\begin{tabular}{@{}l@{\ \ }l@{\qquad}l@{\ \ }l@{\qquad}l@{\ \ }l@{}}
11.9 & \cod{debilitamiento} & 11.14 & \cod{pertenencia_literal} & 11.17 & \cod{descarga}\\
11.10 & \cod{introduccion_doble_negacion} & 11.15 & \cod{kalmar} & 11.18 & \cod{completitud_ND}\\
11.11 & \cod{analisis_por_casos} & 11.16 & \cod{congruencia_contextoLiterales} & A.3, A.6 & \cod{correccion}, \cod{correccion_completitud}\\
\end{tabular}
\end{tcolorbox}

\navegacion{nivel14}{Nivel 14}{nivel16}{Nivel 16}

\newpage
%=====================================================================
% NIVEL 16
%=====================================================================
\nivel{16}{4}{\faCheckDouble}{Completitud: el teorema, línea a línea}{poder
explicar las 14 líneas del teorema principal de la tesis.}

\begin{minipage}[t]{0.60\textwidth}
\vspace{0pt}
\begin{leanS}[title={Thesis/Prop/Completeness.lean (completo)}]
theorem completitud_ND (φ : Formula) (h : EsTautologia φ) :
    ND (∅ : Set Formula) φ := by
  let listaAtomos := (atomos φ).dedup                              @1
  have hSinRepeticiones : listaAtomos.Nodup := List.nodup_dedup _  @2
  have hCobertura : atomos φ ⊆ listaAtomos := by                   @3
    intro a ha
    exact List.mem_dedup.mpr ha
  have hUniversal : ∀ v, ND (contextoLiterales v listaAtomos) φ := by   @4
    intro v
    have hVerdadero : evaluar v φ := h v                           @5
    have hPositiva : evaluar v φ → ND (contextoLiterales v listaAtomos) φ :=
      (kalmar v listaAtomos φ hCobertura).left                     @6
    exact hPositiva hVerdadero
  exact descarga listaAtomos hSinRepeticiones hUniversal           @7
\end{leanS}
\begin{estado}[title={Estado real justo antes de la última línea}]
φ : Formula
h : EsTautologia φ
listaAtomos : List String := (atomos φ).dedup
hSinRepeticiones : listaAtomos.Nodup
hCobertura : atomos φ ⊆ listaAtomos
hUniversal : ∀ (v : Valuacion), ND (contextoLiterales v listaAtomos) φ
⊢ ND ∅ φ
\end{estado}
\end{minipage}\hfill
\begin{minipage}[t]{0.37\textwidth}
\vspace{0pt}
\ex{1}{La lista de átomos de $\varphi$, sin repetidos.}
\ex{2}{Sin repetidos: lo pide \cod{descarga}.}
\ex{3}{Cubre todos los átomos: lo pide \cod{kalmar}.}
\ex{4}{«Para toda fila $v$, sus literales derivan $\varphi$.» Es lo que pide \cod{descarga}.}
\ex{5}{$h$ es un $\forall$: aplicarlo a $v$ da que $\varphi$ es verdadera en esa fila.}
\ex{6}{Kalmár da dos mitades; \cod{.left} toma la positiva.}
\ex{7}{\cod{descarga} recibe sus tres ingredientes y devuelve $\varnothing\vdash\varphi$.}
\clave{m4a}{\textbf{El estado lo dice todo:} tres hechos arriba, que son justo los tres
argumentos de \cod{descarga}.}
\dice{«Tomo los átomos sin repetir. Por Kalmár, cada fila deriva $\varphi$ desde sus
literales, porque $\varphi$ es tautología. La descarga elimina los literales.»}
\end{minipage}

\vspace{1pt}
\begin{lean}[title={Y con la corrección: el «si y solo si» (Teorema A.6; líneas compactadas con «;»)}]
theorem correccion_completitud (φ : Formula) : EsTautologia φ ↔ ND (∅ : Set Formula) φ := by
  constructor                                  -- ↔ : dos direcciones
  · intro hTautologia; exact completitud_ND φ hTautologia        -- ida: completitud
  · intro d; exact tautologia_de_derivable d                     -- vuelta: corrección
\end{lean}

\navegacion{nivel15}{Nivel 15}{nivel17}{Nivel 17}

\newpage
%=====================================================================
% NIVEL 17
%=====================================================================
\nivel{17}{4}{\faExchange*}{Quicksort: la permutación}{leer la prueba de que
quicksort no pierde ni inventa elementos: cuatro hechos y una cadena.}

\begin{minipage}[t]{0.60\textwidth}
\vspace{0pt}
\begin{leanS}[title={quicksort\_permutacion, caso p :: cola (segunda mitad; sangría reducida)}]
have hDerecha : p :: quicksort mayores ~ p :: mayores :=
  hipIndMayores.cons p                                              @1
have hConcatenacion : quicksort menoresIguales ++ (p :: quicksort mayores)
    ~ menoresIguales ++ (p :: mayores) :=
  hipIndMenoresIguales.append hDerecha                              @2
have hPivote : menoresIguales ++ (p :: mayores)
    ~ p :: (menoresIguales ++ mayores) := List.perm_middle          @3
have hCons : p :: (menoresIguales ++ mayores) ~ p :: cola :=
  hParticion.cons p                                                 @4
rw [quicksort_cons]                                                 @5
change quicksort menoresIguales ++ (p :: quicksort mayores) ~ p :: cola
calc                                                                @6
  quicksort menoresIguales ++ (p :: quicksort mayores)
      ~ menoresIguales ++ (p :: mayores) := hConcatenacion
  _ ~ p :: (menoresIguales ++ mayores) := hPivote
  _ ~ p :: cola := hCons
\end{leanS}
\begin{tcolorbox}[colback=m4b!7, colframe=m4b, boxrule=1pt, arc=4pt, left=6pt, right=6pt,
  top=3pt, bottom=3pt, halign=center]
\textbf{La misma cadena en papel} \ ($L$: menores o iguales, $R$: mayores)\\[3pt]
$\mathrm{qs}(L)\mathbin{+\!\!+}(p::\mathrm{qs}(R))\ \overset{\text{\bola{2}}}{\sim}\ L\mathbin{+\!\!+}(p::R)
\ \overset{\text{\bola{3}}}{\sim}\ p::(L\mathbin{+\!\!+}R)\ \overset{\text{\bola{4}}}{\sim}\ p::\mathrm{cola}$
\end{tcolorbox}
\end{minipage}\hfill
\begin{minipage}[t]{0.37\textwidth}
\vspace{0pt}
\ex{1}{\cod{.cons p}: añadir el mismo elemento delante conserva la permutación.}
\ex{2}{\cod{.append}: pegar dos permutaciones da una permutación. Usa las \textbf{dos hipótesis inductivas}.}
\ex{3}{\cod{List.perm_middle} (Mathlib): mover $p$ del medio al frente.}
\ex{4}{\cod{hParticion}: los dos filtros juntos son una permutación de \cod{cola} (Lema 12.7).}
\ex{5}{Despliega \cod{quicksort (p :: cola)} con su ecuación.}
\ex{6}{\cod{calc} encadena por transitividad.}
\clave{m4b}{\textbf{Ojo:} aquí \cod{~} es «\textbf{es permutación de}» (\cod{List.Perm}), no la negación de fórmulas del capítulo 11.}
\tarj{m4b}{Qué viene de Mathlib}{La relación \cod{~} y sus reglas básicas (\cod{cons}, \cod{append}, \cod{perm_middle}). \textbf{Ningún} teorema sobre quicksort.}
\end{minipage}

\pruebate{m4b}{\item ¿Dónde se usan las hipótesis inductivas?
\item ¿Qué lema propio justifica el último eslabón?
\item ¿Qué hace \cod{hipIndMayores.cons p}?}
{1.\ En \bola{1} y \bola{2}: una por cada mitad.
2.\ \cod{particion_filtros_permutacion} (Lema 12.7), a través de \cod{hParticion}.
3.\ De \cod{quicksort mayores ~ mayores} obtiene lo mismo con $p$ delante de los dos lados.}

\navegacion{nivel16}{Nivel 16}{nivel18}{Nivel 18}

\newpage
%=====================================================================
% NIVEL 18
%=====================================================================
\nivel{18}{4}{\faSortAmountUp}{Quicksort: el orden y el teorema final}{leer por
qué la salida queda ordenada y cómo se juntan las dos mitades de la especificación.}

\begin{minipage}[t]{0.60\textwidth}
\vspace{0pt}
\begin{leanS}[title={quicksort\_ordenada, caso p :: cola (los cinco hechos; las pruebas «…» están recortadas)}]
have hCotaMenoresIguales : ∀ a ∈ quicksort menoresIguales, a ≤ p := …   @1
have hCotaMayores : ∀ b ∈ quicksort mayores, p ≤ b := …                 @2
have hParesDerecha : List.Pairwise (· ≤ ·) (p :: quicksort mayores) := …  @3
have hCruce : ∀ a ∈ quicksort menoresIguales,
    ∀ b ∈ p :: quicksort mayores, a ≤ b := …                            @4
have hTodosPares : List.Pairwise (· ≤ ·)
    (quicksort menoresIguales ++ (p :: quicksort mayores)) := by         @5
  apply (relacion_pares_concatenacion_iff _ _ _).mpr
  exact ⟨hParesMenoresIguales, hParesDerecha, hCruce⟩
\end{leanS}
\begin{lean}[title={El teorema final (Teorema 12.12), completo}]
theorem quicksort_correcto (l : List α) :
    quicksort l ~ l ∧ Ordenada (quicksort l) := by
  have hPermutacion : quicksort l ~ l := quicksort_permutacion l
  have hOrdenada : Ordenada (quicksort l) := quicksort_ordenada l
  exact And.intro hPermutacion hOrdenada
\end{lean}
\end{minipage}\hfill
\begin{minipage}[t]{0.37\textwidth}
\vspace{0pt}
\ex{1}{Todo lo de la izquierda es $\le p$: estaba en el filtro (por la permutación ya probada).}
\ex{2}{Todo lo de la derecha es $\ge p$. De \cod{¬ (b ≤ p)} sale \cod{p < b}: aquí se usa que el orden es \textbf{lineal}.}
\ex{3}{$p$ delante de la derecha ordenada sigue ordenado.}
\ex{4}{Cada elemento izquierdo es $\le$ cada elemento del bloque derecho (transitividad).}
\ex{5}{El lema de concatenación (12.10), en la dirección \cod{.mpr}: de las tres condiciones a \cod{Pairwise}.}
\clave{m4b}{\cod{Ordenada l} es \cod{l.Pairwise (· ≤ ·)}: todo elemento es $\le$ que cada uno de los posteriores.}
\clave{m4b}{El final es una línea: \cod{And.intro} junta las dos pruebas. Es el \cod{∧} de la tabla maestra \chip{8}.}
\end{minipage}

\vspace{2pt}
\begin{tcolorbox}[colback=m4b!6, colframe=m4b, boxrule=1pt, arc=4pt, left=6pt, right=6pt,
  top=3pt, bottom=3pt, title={\faExchange*\ Del papel a Lean: número del resultado y nombre}, colbacktitle=m4b,
  fonttitle=\bfseries]
\small
\begin{tabular}{@{}l@{\ \ }l@{\qquad}l@{\ \ }l@{\qquad}l@{\ \ }l@{}}
12.5 & \cod{longitud_filtro_le} & 12.8 & \cod{quicksort_permutacion} & 12.11 & \cod{quicksort_ordenada}\\
12.6 & \cod{termination_by}, \cod{longitud_filtro_lt_cons} & 12.9 & \cod{pertenencia_quicksort} & 12.12 & \cod{quicksort_correcto}\\
12.7 & \cod{particion_filtros_permutacion} & 12.10 & \cod{relacion_pares_concatenacion_iff} & & \\
\end{tabular}
\end{tcolorbox}

\navegacion{nivel17}{Nivel 17}{nivel19}{Nivel 19}

\newpage
%=====================================================================
% NIVEL 19
%=====================================================================
\nivel{19}{4}{\faShield*}{Auditoría: qué comprobó Lean}{poder decir con
exactitud de qué depende cada teorema y qué garantiza la verificación.}

\begin{minipage}[t]{0.52\textwidth}
\vspace{0pt}
\begin{estado}[title={Salida real de \#print axioms (reports/*/audit.txt)}]
'Thesis.Prop.completitud_ND' depends on axioms:
    [propext, Classical.choice, Quot.sound]
'Thesis.Prop.correccion_completitud' depends on axioms:
    [propext, Classical.choice, Quot.sound]
'Thesis.Sort.quicksort_correcto' depends on axioms:
    [propext, Quot.sound]
'Thesis.Prop.ej_identidad' does not depend on any axioms
\end{estado}
\begin{lean}[title={Si hubiera un hueco, se vería (ejemplo aparte, no es de la tesis)}]
theorem falso_admitido : 1 = 2 := by sorry
#print axioms falso_admitido
-- 'falso_admitido' depends on axioms: [sorryAx]
\end{lean}
\begin{tcolorbox}[colback=m4c!6, colframe=m4c, boxrule=1pt, arc=4pt, left=6pt, right=6pt,
  top=3pt, bottom=3pt, title={\faTerminal\ Cómo se repite}, colbacktitle=m4c, fonttitle=\bfseries]
\cod{lake build} \hfill compila todo: si termina sin errores, el núcleo aceptó cada prueba\\
\cod{lake env lean Thesis/Prop/Audit.lean} \hfill imprime los axiomas\\
\cod{make audit} \hfill regenera el certificado; la integración continua lo compara
\end{tcolorbox}
\end{minipage}\hfill
\begin{minipage}[t]{0.45\textwidth}
\vspace{0pt}
\begin{tcolorbox}[colback=m4c!6, colframe=m4c, boxrule=1pt, arc=4pt, left=6pt, right=6pt,
  top=3pt, bottom=3pt, title={\faKey\ Los tres axiomas}, colbacktitle=m4c, fonttitle=\bfseries]
\renewcommand{\arraystretch}{1.2}
\begin{tabular}{@{}l@{\quad}>{\raggedright\arraybackslash}p{6.9cm}@{}}
\cod{propext} & proposiciones equivalentes son iguales\\
\cod{Quot.sound} & elementos relacionados dan la misma clase en un cociente\\
\cod{Classical.choice} & elegir un elemento de un tipo no vacío; da la lógica clásica\\
\end{tabular}
\end{tcolorbox}
\clave{m4c}{\textbf{Quicksort no usa \cod{Classical.choice}.} Completitud sí: razona por
casos y por absurdo sobre \cod{evaluar v φ}, que es una \cod{Prop} cualquiera \chip{10}.}
\begin{tcolorbox}[colback=white, colframe=black!75, boxrule=1.3pt, arc=6pt,
  title={\faBalanceScale\ Qué garantiza Lean y qué no}, colbacktitle=black!75,
  fonttitle=\bfseries, left=6pt, right=6pt, top=3pt, bottom=3pt]
{\color{verde}\faCheckCircle}\ Que cada prueba es un término del tipo de su enunciado.\\
{\color{verde}\faCheckCircle}\ Que no hay \cod{sorry} ni axiomas fuera de la lista.\\[3pt]
{\color{rojo}\faTimesCircle}\ Que el enunciado diga lo que tú crees.\\
{\color{rojo}\faTimesCircle}\ Que las definiciones (\cod{ND}, \cod{evaluar}, \cod{Ordenada}) sean las correctas.\\[3pt]
Por eso los enunciados son cortos y se comparan con el papel.
\end{tcolorbox}
\end{minipage}

\pruebate{m4c}{\item ¿De qué axiomas depende \cod{quicksort_correcto}?
\item ¿Qué pasaría en la auditoría si quedara un \cod{sorry}?
\item ¿Por qué \cod{ej_identidad} no depende de ningún axioma?}
{1.\ De \cod{propext} y \cod{Quot.sound}; no de \cod{Classical.choice}.
2.\ Aparecería \cod{sorryAx} en la lista.
3.\ Es una derivación construida solo con constructores de \cod{ND}: datos puros.}

\navegacion{nivel18}{Nivel 18}{simbolos}{Símbolos}

\newpage
%=====================================================================
% SÍMBOLOS
%=====================================================================
\pagina{simbolos}{m2}{\faKeyboard}{Chuleta de símbolos: cómo se leen}{si un
símbolo te frena, está aquí. La tercera columna es cómo se teclea en VS Code.}

\newcommand{\sy}[3]{\cod{#1} & #2 & {\color{gris}\ttfamily\small #3}\\}
\begin{minipage}[t]{0.325\textwidth}
\vspace{0pt}
\begin{tcolorbox}[colback=m2!5, colframe=m2, boxrule=1pt, arc=5pt, left=5pt, right=5pt,
  title={Lógica (de Lean)}, colbacktitle=m2, fonttitle=\bfseries]
\renewcommand{\arraystretch}{1.3}
\begin{tabular}{@{}>{\large}l@{\ \ }>{\raggedright}p{3.85cm}@{\ }l@{}}
\sy{→}{implica; función}{\textbackslash to}
\sy{∀}{para todo}{\textbackslash forall}
\sy{∃}{existe}{\textbackslash exists}
\sy{¬}{no}{\textbackslash not}
\sy{∧}{y}{\textbackslash and}
\sy{∨}{o}{\textbackslash or}
\sy{↔}{si y solo si}{\textbackslash iff}
\sy{=}{es igual a (una \cod{Prop})}{}
\sy{:=}{se define como}{}
\sy{⊢}{«hay que probar» (solo en el estado)}{}
\end{tabular}
\end{tcolorbox}
\end{minipage}\hfill
\begin{minipage}[t]{0.325\textwidth}
\vspace{0pt}
\begin{tcolorbox}[colback=m4a!5, colframe=m4a, boxrule=1pt, arc=5pt, left=5pt, right=5pt,
  title={Datos y notación de la tesis}, colbacktitle=m4a, fonttitle=\bfseries]
\renewcommand{\arraystretch}{1.3}
\begin{tabular}{@{}>{\large}l@{\ \ }>{\raggedright}p{3.85cm}@{\ }l@{}}
\sy{~φ}{negación de una fórmula}{}
\sy{φ ⟶ ψ}{implicación de fórmulas}{\textbackslash hom}
\sy{l₁ ~ l₂}{es permutación de}{}
\sy{∈}{pertenece a}{\textbackslash in}
\sy{⊆}{está contenido en}{\textbackslash sub}
\sy{∅}{vacío}{\textbackslash empty}
\sy{::}{poner delante de la lista}{}
\sy{++}{pegar dos listas}{}
\sy{≤}{menor o igual}{\textbackslash le}
\sy{!}{«no» de un \cod{Bool}}{}
\end{tabular}
\end{tcolorbox}
\end{minipage}\hfill
\begin{minipage}[t]{0.325\textwidth}
\vspace{0pt}
\begin{tcolorbox}[colback=m3!5, colframe=m3, boxrule=1pt, arc=5pt, left=5pt, right=5pt,
  title={Escritura de pruebas}, colbacktitle=m3, fonttitle=\bfseries]
\renewcommand{\arraystretch}{1.3}
\begin{tabular}{@{}>{\large}l@{\ \ }>{\raggedright}p{3.85cm}@{\ }l@{}}
\sy{⟨a, b⟩}{par: constructor anónimo}{\textbackslash< \textbackslash>}
\sy{·}{enfoca una meta}{\textbackslash .}
\sy{_}{dedúcelo tú}{}
\sy{?_}{meta que pruebo después}{}
\sy{fun x => e}{la función $x\mapsto e$}{}
\sy{(· ≤ ·)}{\cod{fun a b => a ≤ b}}{}
\sy{h ▸ e}{reescribe \cod{e} con \cod{h}}{\textbackslash t}
\sy{@f}{\cod{f} con todo explícito}{}
\sy{x✝}{nombre inaccesible}{}
\sy{--}{comentario}{}
\end{tabular}
\end{tcolorbox}
\end{minipage}

\vspace{4pt}
\begin{tcbraster}[raster columns=3, raster column skip=6pt, raster equal height]
\clave{m2}{\textbf{Letras griegas:} \cod{Γ}, \cod{Δ} son contextos; \cod{φ}, \cod{ψ}, \cod{χ} son fórmulas; \cod{α} es un tipo cualquiera. Se teclean \texttt{\textbackslash Gamma}, \texttt{\textbackslash phi}…}
\clave{m4a}{\textbf{Subíndices:} \cod{l₁} es solo un nombre, como \cod{l1}. Se teclea \texttt{\textbackslash 1}.}
\clave{m3}{\textbf{El punto:} \cod{h.left}, \cod{h.mp}, \cod{l.length}: «de \cod{h}, su parte…». \cod{.atom s} al inicio: el constructor, con el tipo sobreentendido.}
\end{tcbraster}

\navegacion{nivel19}{Nivel 19}{diccionario}{Cómo decirlo}

\newpage
%=====================================================================
% DICCIONARIO
%=====================================================================
\pagina{diccionario}{m3}{\faCommentDots}{Cómo decirlo en voz alta, y 10 trampas}{el
tribunal evalúa el argumento, no la sintaxis. Traduce cada línea a una frase.}

\begin{minipage}[t]{0.50\textwidth}
\vspace{0pt}
\begin{tcolorbox}[colback=white, colframe=m3, boxrule=1.1pt, arc=5pt, left=5pt, right=5pt,
  top=2pt, bottom=2pt, title={\faCommentDots\ De Lean al español}, colbacktitle=m3, fonttitle=\bfseries\large]
\renewcommand{\arraystretch}{1.24}
\begin{tabular}{@{}l@{\quad}>{\raggedright\arraybackslash}p{7.6cm}@{}}
\cod{intro h} & «Supongamos…» / «Sea … cualquiera.»\\
\cod{have h : P := …} & «Primero observamos que $P$.»\\
\cod{exact e} & «…que es justo lo que había que probar.»\\
\cod{apply f} & «Por $f$, basta probar…»\\
\cod{constructor} & «Probamos cada parte por separado.»\\
\cod{obtain ⟨a, b⟩ := h} & «De $h$ obtenemos $a$ y $b$.»\\
\cod{rcases h with a | b} & «Distinguimos dos casos.»\\
\cod{by_cases hp : P} & «O bien $P$, o bien no $P$.»\\
\cod{by_contra h} & «Por reducción al absurdo, supongamos que no.»\\
\cod{induction x with} & «Por inducción sobre $x$.»\\
\cod{rw [h]} & «Reescribiendo con $h$…»\\
\cod{simp only [...]} & «Simplificando con estos lemas…»\\
\cod{change P} & «Por definición, la meta es $P$.»\\
\cod{calc} & «Encadenando: … que a su vez …»\\
\cod{hipInd hAmpliado} & «Por la hipótesis inductiva, aplicada a…»\\
\cod{ND.impI dCuerpo} & «Por la regla $\to$I, descargando la hipótesis.»\\
\end{tabular}
\end{tcolorbox}
\end{minipage}\hfill
\begin{minipage}[t]{0.47\textwidth}
\vspace{0pt}
\begin{tcolorbox}[colback=white, colframe=rojo, boxrule=1.1pt, arc=5pt, left=5pt, right=5pt,
  top=2pt, bottom=2pt, title={\faExclamationTriangle\ 10 trampas}, colbacktitle=rojo, fonttitle=\bfseries\large]
\begin{enumerate}[leftmargin=6mm, itemsep=1.5pt, topsep=1pt]
\item \cod{~} significa dos cosas: negación de fórmulas (cap.\ 11) y «es permutación de» (cap.\ 12).
\item \cod{→} es a la vez «implica» y «función». Por Curry–Howard es lo mismo.
\item \cod{¬}, \cod{→}, \cod{False} son de Lean. \cod{~}, \cod{⟶}, \cod{Falsedad} son del cálculo.
\item \cod{hipInd}, \cod{hsub}, \cod{dP} son solo nombres. Lean no les da significado.
\item \cod{_} es «dedúcelo tú»; \cod{?_} es «lo pruebo después».
\item \cod{:=} define; \cod{=} es la igualdad, una proposición.
\item \cod{Set} (sin orden ni repeticiones) no es \cod{List} (con orden y repeticiones).
\item \cod{decide (x ≤ p)} es un \cod{Bool}; \cod{x ≤ p} es una \cod{Prop}.
\item \cod{change} no demuestra nada: solo muestra la misma meta desplegada.
\item \cod{ND.classical} es un dato; \cod{by_contra} es lógica clásica de Lean.
\end{enumerate}
\end{tcolorbox}
\clave{m3}{\large\textbf{Receta para explicar una prueba:} (1) el enunciado en español;
(2) el método, del comentario \cod{-- Método:}; (3) el estado en el paso clave;
(4) qué lema lo cierra.}
\end{minipage}

\navegacion{simbolos}{Símbolos}{indice}{Índice}

\newpage
%=====================================================================
% ÍNDICE
%=====================================================================
\pagina{indice}{m4c}{\faFont}{Índice: de la palabra al nivel}{pulsa el número
para ir al nivel donde se explica.}

\newcommand{\ix}[2]{\cod{#1}\ \chip{#2}\par}
\newcommand{\ixx}[3]{\cod{#1}\ \chip{#2}\ \chip{#3}\par}
\begin{tcolorbox}[colback=white, colframe=m4c, boxrule=1pt, arc=6pt, top=3pt, bottom=3pt]
\begin{minipage}[t]{0.19\textwidth}\raggedright
\texttt{\#check}\ \chip{1}\par
\texttt{\#eval}\ \chip{1}\par
\texttt{\#print axioms}\ \chip{19}\par
\ix{abbrev}{2}\ix{And.intro}{18}\ixx{apply}{7}{8}\ix{at}{9}\ix{by}{5}\ix{by_cases}{10}\ix{by_contra}{10}\ix{calc}{9}\ix{cases … with}{10}
\end{minipage}\hfill
\begin{minipage}[t]{0.19\textwidth}\raggedright
\ix{change}{9}\ix{Classical.choice}{19}\ixx{constructor}{7}{8}\ix{decide}{13}\ix{decreasing_by}{13}\ix{def}{4}\ix{deriving}{2}\ix{exact}{7}\ix{example}{5}\ix{False.elim}{8}\ix{fun … => …}{4}\ix{have}{7}
\end{minipage}\hfill
\begin{minipage}[t]{0.19\textwidth}\raggedright
\ix{hipInd}{11}\ix{import}{2}\ix{induction … with}{11}\ix{inductive}{3}\ix{insert φ Γ}{3}\ix{intro}{7}\ix{lake build}{19}\ix{let}{14}\ix{List}{3}\ix{.mp / .mpr}{8}\ix{namespace}{2}\ix{ND}{12}
\end{minipage}\hfill
\begin{minipage}[t]{0.19\textwidth}\raggedright
\ix{noncomputable}{4}\ix{obtain}{8}\ix{open}{2}\ix{Prop / Type}{1}\ix{rcases}{10}\ix{refine ⟨?_, ?_⟩}{7}\ix{rfl}{9}\ix{rw}{9}\ix{Set}{3}\ix{simp only}{9}\ix{simpa}{9}\ix{sorry}{19}
\end{minipage}\hfill
\begin{minipage}[t]{0.19\textwidth}\raggedright
\ix{subst}{9}\ix{termination_by}{13}\ix{theorem}{5}\ix{variable}{4}\ix{⊢}{6}\ix{⟨ , ⟩}{8}\ix{·}{7}\ix{_}{4}\ix{@}{12}\ix{{Γ φ}}{4}\ix{[LinearOrder α]}{4}\ix{~}{17}
\end{minipage}
\end{tcolorbox}

\newcommand{\grupo}[3]{\begin{tcolorbox}[colback=#1!6, colframe=#1, boxrule=1.1pt,
  arc=5pt, title={#2}, colbacktitle=#1, fonttitle=\bfseries, left=5pt, top=2pt, bottom=2pt]
  \begin{itemize}[label={\faSquare[regular]}, leftmargin=6mm, itemsep=1.5pt, topsep=0pt]#3\end{itemize}\end{tcolorbox}}
{\large\bfseries\faTasks\ ¿Listo? Marca cuando lo hagas sin mirar}\\[1pt]
\begin{tcbraster}[raster columns=4, raster column skip=6pt, raster equal height]
\grupo{m1}{Mundo 1 · Leer}{
  \item Leer \cod{a : A} y un \texttt{\#check}
  \item Decir un enunciado en español
  \item Pasar $\Gamma,\varphi\vdash\psi$ a Lean}
\grupo{m2}{Mundo 2 · Probar}{
  \item Leer un estado: qué tengo, qué falta
  \item La tabla maestra de memoria
  \item Objeto y meta: tres parejas}
\grupo{m3}{Mundo 3 · Inducción}{
  \item Qué dice \cod{hipInd} en cada caso
  \item Las 7 reglas $=$ 7 constructores
  \item Por qué \cod{termination_by}}
\grupo{m4a}{Mundo 4 · Teoremas}{
  \item \cod{completitud_ND} en 7 pasos
  \item La cadena \cod{calc} de quicksort
  \item Los axiomas de cada teorema}
\end{tcbraster}

\navegacion{diccionario}{Cómo decirlo}{portada}{Portada}

\end{document}
